Esta nota describe las fuentes, las transformaciones, los métodos
estadísticos, las reglas de clasificación y los controles de calidad que
producen cada cifra del sitio ColombiaMacro. Reemplaza a la versión 9
(septiembre de 2026), que cubría solo las secciones originales del
tablero. Todo lo que aquí se describe corresponde al código vigente del
repositorio; cuando una afirmación depende de un resultado numérico, se
presenta como ejercicio fechado (sección 8) y no como propiedad
permanente del método.

\hypertarget{propuxf3sito-preguntas-y-principios-de-diseuxf1o}{%
\subsection{1. Propósito, preguntas y principios de
diseño}\label{propuxf3sito-preguntas-y-principios-de-diseuxf1o}}

\hypertarget{propuxf3sito}{%
\subsubsection{1.1 Propósito}\label{propuxf3sito}}

ColombiaMacro es un tablero académico, bilingüe (español e inglés),
dirigido a inversionistas, analistas y estudiantes que necesitan una
lectura ordenada de la economía colombiana a partir de datos oficiales.
Su pregunta central es \textbf{en qué fase del ciclo está la economía
colombiana y qué implica para los precios, la política monetaria, las
cuentas externas y fiscales y los activos}. Cada página responde una
pregunta específica y abre con un bloque de ``lo clave'' cuyas frases se
generan automáticamente a partir de los datos.

\hypertarget{puxe1ginas-y-preguntas}{%
\subsubsection{1.2 Páginas y preguntas}\label{puxe1ginas-y-preguntas}}

\small\begin{longtable}[]{@{}
  >{\raggedright\arraybackslash}p{(\columnwidth - 4\tabcolsep) * \real{0.3333}}
  >{\raggedright\arraybackslash}p{(\columnwidth - 4\tabcolsep) * \real{0.3333}}
  >{\raggedright\arraybackslash}p{(\columnwidth - 4\tabcolsep) * \real{0.3333}}@{}}
\toprule\noalign{}
\begin{minipage}[b]{\linewidth}\raggedright
Página
\end{minipage} & \begin{minipage}[b]{\linewidth}\raggedright
Pregunta que responde
\end{minipage} & \begin{minipage}[b]{\linewidth}\raggedright
Medidas principales
\end{minipage} \\
\midrule\noalign{}
\endhead
\bottomrule\noalign{}
\endlastfoot
Portada & ¿Cómo está la economía hoy, en un vistazo? & Fase del ciclo,
seis tarjetas con estado, curva TES, resumen por tema \\
Ciclo & ¿En qué parte del ciclo está la economía? & Brecha del producto
(cinco métodos), reloj del ciclo, ciclo mensual del ISE, giros, Okun \\
Crecimiento & ¿Cuánto crece la economía y por qué? & Tres velocidades
del PIB, deflactor, aportes por sector y por componente del gasto,
inversión, consumo, PIB por persona, productividad \\
Capacidad & ¿Produce por encima o por debajo de su capacidad, y dónde? &
Brecha agregada y sectorial, holgura laboral, regiones (departamentos),
diversificación \\
Empleo & ¿Cómo está el empleo y cuánto es informal? & Desempleo,
informalidad, empleo por rama y posición, brechas por sexo, jóvenes,
zonas, 32 capitales \\
Inflación & ¿Converge la inflación a la meta y de dónde viene? & IPC
total y de fondo, breakevens, forward 5y5y, divisiones, bienes y
servicios, difusión, ingresos, ciudades \\
Tasas & ¿Qué hace el Banco de la República y cómo se transmite? & Tasa
de política, tasa real ex ante, ciclos, traspaso, crédito, IBR,
liquidez \\
Curva TES & ¿Cuánto cobra el mercado por prestarle al Gobierno? & Curva
cero cupón, nivel-pendiente-curvatura, episodios invertidos, Fisher,
prima, volatilidad \\
Mercados & ¿Cómo está el peso y qué lo mueve? & TRM, pares regionales,
dólar global, ITCR, petróleo, balanza cambiaria, reservas,
volatilidad \\
Empresas & ¿Cómo están las empresas y qué mueve la bolsa? & COLCAP,
bolsa por dentro, 10.000 empresas, crédito empresarial, registro
mercantil \\
Externo & ¿Cómo están las cuentas con el exterior y del Gobierno? &
Balanza de pagos (MBP6), financiación, IED, remesas, deuda externa, PII,
balance fiscal \\
Comercio & ¿Qué vende y qué compra Colombia, y con quién? &
Exportaciones e importaciones, CUODE, precio y volumen, socios,
apertura, concentración \\
Indicadores / datos & ¿Qué tan extremo es cada dato y de dónde viene? &
16 indicadores con percentil histórico, publicaciones oficiales, estado
de las fuentes, descargas \\
\end{longtable}\normalsize

\hypertarget{principios-de-diseuxf1o}{%
\subsubsection{1.3 Principios de diseño}\label{principios-de-diseuxf1o}}

\begin{enumerate}
\def\labelenumi{\arabic{enumi}.}
\tightlist
\item
  \textbf{Fuentes oficiales.} Las cifras provienen del DANE, el Banco de
  la República (BanRep), el Ministerio de Hacienda (vía BanRep), la
  Superintendencia Financiera (vía BanRep), la Superintendencia de
  Sociedades y Confecámaras (vía el portal de datos abiertos
  datos.gov.co) y, para el contexto global, la Reserva Federal de EE.
  UU. y la Administración de Información Energética (EIA) vía FRED. Las
  únicas excepciones, explícitamente señaladas como fuentes auxiliares
  no oficiales, son la composición del fondo iShares MSCI COLCAP y los
  precios de Yahoo Finance usados en ``la bolsa por dentro'' (sección
  3.4).
\item
  \textbf{Descarga automática y verificable.} Ninguna cifra se digita a
  mano, con una sola excepción documentada: el archivo de decisiones
  anunciadas por la Junta del Banco de la República (sección 3.5). Cada
  serie del graficador del Banco se acepta solo si su identificador
  \textbf{y} su nombre coinciden con lo esperado.
\item
  \textbf{Sin proyecciones ni pronósticos.} El sitio describe el
  presente con datos observados. Las trayectorias ``hacia el futuro''
  que aparecen (por ejemplo, la inflación implícita en los TES) son
  precios de mercado observados hoy, no pronósticos del tablero.
\item
  \textbf{Sin información futura en los cálculos ``en tiempo real''.} La
  brecha principal se estima con un filtro de una cola que, para cada
  trimestre, solo usa datos hasta ese trimestre; la tasa real ex post
  usa la inflación ya publicada en cada fecha. Las medidas que sí usan
  toda la muestra (filtros de dos colas, coeficientes estimados con la
  muestra completa) se identifican como tales.
\item
  \textbf{Incertidumbre visible.} La brecha del producto se presenta con
  cinco métodos, su mediana y su rango; las inflaciones implícitas se
  rotulan como compensación por inflación que incluye primas, no como
  expectativas puras.
\item
  \textbf{Reglas explícitas.} Cada estado (``expansión'',
  ``restrictiva'', ``por encima de la meta'') y cada frase del veredicto
  se deriva de un umbral escrito en el código (sección 6). Los estados
  describen; no son recomendaciones de inversión.
\item
  \textbf{Un eje, una unidad.} Los gráficos principales usan un solo
  eje; cuando se combinan magnitudes distintas (por ejemplo, saldo en
  dólares y razón al PIB) el eje secundario se rotula explícitamente.
\end{enumerate}

\hypertarget{convenciones-y-notaciuxf3n}{%
\subsection{2. Convenciones y
notación}\label{convenciones-y-notaciuxf3n}}

\begin{itemize}
\tightlist
\item
  \textbf{Fechas.} Los datos mensuales se fechan el primer día del mes;
  los trimestrales, el primer día del trimestre; los anuales, el 1 de
  enero. Los trimestres móviles y años móviles de la GEIH se fechan en
  su último mes. Las series semanales de acciones se fechan el viernes
  de la semana (o el día del último dato si la semana está en curso).
\item
  \textbf{Tasas.} Todas las tasas de interés están en porcentaje
  efectivo anual. Las conversiones entre nominal y real usan la ecuación
  de Fisher exacta, \((1+i) = (1+r)(1+\pi)\), no la aproximación
  \(i-\pi\).
\item
  \textbf{Variaciones.} Para una serie \(x\), la variación anual es
  \(100\,(x_t/x_{t-k}-1)\) con \(k=12\) en datos mensuales, \(k=4\) en
  trimestrales y \(k=52\) en semanales. Los cambios de tasas se expresan
  en puntos porcentuales (pp) o puntos básicos (pb, 1 pp = 100 pb).
\item
  \textbf{``Último dato disponible''.} Cuando una comparación se hace
  ``frente a hace un año'', se toma el último dato disponible en o antes
  de la fecha de referencia menos el horizonte (365, 364, 91 o 28 días
  según la medida). En los paneles de cambio de los gráficos se busca el
  dato más cercano con una tolerancia de 20 días (horizontes anual y
  trimestral) o 12 días (horizonte mensual).
\item
  \textbf{Brechas en logaritmos.} Con \(y_t = 100\ln Y_t\) y una
  tendencia \(\tau_t\) en la misma escala, la brecha
  \(g_t = y_t - \tau_t\) se lee como porcentaje aproximado del nivel de
  tendencia.
\item
  \textbf{Unidades monetarias.} ``Billones'' designa \(10^{12}\) pesos;
  ``miles de millones de dólares'', \(10^{9}\) dólares. Los niveles
  reales del PIB están en volúmenes encadenados con referencia 2015.
\end{itemize}

\hypertarget{fuentes-de-datos}{%
\subsection{3. Fuentes de datos}\label{fuentes-de-datos}}

\hypertarget{tabla-de-archivos}{%
\subsubsection{3.1 Tabla de archivos}\label{tabla-de-archivos}}

La tabla enumera cada archivo de \texttt{data/} que usa el sitio.
``Frecuencia'' es la de la serie; ``Identificador'' es el número de
serie del graficador SUAMECA del Banco de la República, el identificador
del conjunto de datos de datos.gov.co, el código FRED o el anexo del
DANE que lee el código.

\small\begin{longtable}[]{@{}
  >{\raggedright\arraybackslash}p{(\columnwidth - 8\tabcolsep) * \real{0.2000}}
  >{\raggedright\arraybackslash}p{(\columnwidth - 8\tabcolsep) * \real{0.2000}}
  >{\raggedright\arraybackslash}p{(\columnwidth - 8\tabcolsep) * \real{0.2000}}
  >{\raggedright\arraybackslash}p{(\columnwidth - 8\tabcolsep) * \real{0.2000}}
  >{\raggedright\arraybackslash}p{(\columnwidth - 8\tabcolsep) * \real{0.2000}}@{}}
\toprule\noalign{}
\begin{minipage}[b]{\linewidth}\raggedright
Archivo
\end{minipage} & \begin{minipage}[b]{\linewidth}\raggedright
Contenido
\end{minipage} & \begin{minipage}[b]{\linewidth}\raggedright
Entidad
\end{minipage} & \begin{minipage}[b]{\linewidth}\raggedright
Identificador
\end{minipage} & \begin{minipage}[b]{\linewidth}\raggedright
Frecuencia
\end{minipage} \\
\midrule\noalign{}
\endhead
\bottomrule\noalign{}
\endlastfoot
\texttt{pib\_colombia.csv} & PIB real original y desestacionalizado
(volúmenes encadenados, ref. 2015) y PIB nominal & DANE & Anexos
\texttt{anex-ProduccionConstantes} (Cuadros 1 y 4) y
\texttt{anex-ProduccionCorriente} (Cuadro 1) & Trimestral \\
\texttt{pib\_sectores.csv} & Valor agregado real de 12 agrupaciones CIIU
Rev.~4: nivel, variación anual, peso y aporte & DANE &
\texttt{anex-ProduccionConstantes}, Cuadro 1 & Trimestral \\
\texttt{pib\_gasto.csv} & PIB por el enfoque del gasto: real original,
real desestacionalizado y nominal & DANE & \texttt{anex-GastoConstantes}
(Cuadros 1, 2, 7 y 8) y \texttt{anex-GastoCorriente} (Cuadro 1) &
Trimestral \\
\texttt{pib\_inversion.csv} & Formación bruta de capital fijo por tipo
de activo (real original y desestacionalizada) & DANE &
\texttt{anex-GastoConstantes}, Cuadros 5 y 6 & Trimestral \\
\texttt{pib\_consumo\_hogares.csv} & Consumo de los hogares por
durabilidad y por finalidad COICOP & DANE &
\texttt{anex-GastoConstantes}, Cuadros 3 y 4 & Trimestral \\
\texttt{poblacion.csv} & Población total nacional a mitad de año
(proyecciones y retroproyecciones, Censo 2018, actualización 2025) &
DANE & \texttt{PPED-AreaNac-2018-20xx} y \texttt{PPED-AreaNac-1950-2017}
& Anual \\
\texttt{pib\_departamentos.csv} & PIB departamental corriente, real y
por habitante & DANE & \texttt{anex-PIBDep-TotalDep} (Cuadros 1 a 3) &
Anual \\
\texttt{pib\_departamentos\_ramas.csv} & Valor agregado corriente por
departamento y 12 ramas, más impuestos & DANE &
\texttt{anex-PIBDep-Activecono}, Cuadro 1 & Anual \\
\texttt{ise\_mensual.csv} & Indicador de Seguimiento a la Economía,
original y desestacionalizado, y tres grandes ramas & DANE &
\texttt{anex-ISE-9actividades} (Cuadros 1 y 2) & Mensual \\
\texttt{mercado\_laboral.csv} & TGP, TO, TD, ocupados, desocupados y
población fuera de la fuerza de trabajo, desestacionalizados & DANE
(GEIH) & \texttt{anex-GEIH-Desestacionalizado}, hoja Total nacional &
Mensual \\
\texttt{laboral\_subutilizacion.csv} & TD, TS y medidas de
subutilización TCSD, TCDFTP y MCSFT (sin desestacionalizar) & DANE
(GEIH) &
\texttt{anex-GEIH-\textless{}mes\textgreater{}\textless{}año\textgreater{}},
hoja Total nacional & Mensual \\
\texttt{laboral\_ciudades.csv} & TD, TS, TGP y TO de las 32 capitales &
DANE (GEIH) &
\texttt{anex-GEIH-\textless{}mes\textgreater{}\textless{}año\textgreater{}},
hoja Año móvil 32 ciudades & Año móvil \\
\texttt{empleo\_ramas.csv} & Ocupados por rama de actividad & DANE
(GEIH) & \texttt{anex-GEIH}, hoja Ocupados TN T13 rama & Mensual \\
\texttt{empleo\_posicion.csv} & Ocupados por posición ocupacional & DANE
(GEIH) & \texttt{anex-GEIH}, hoja Ocupados TN posición & Mensual \\
\texttt{empleo\_sexo.csv} & TGP, TO, TD, TS, ocupados y fuera de la
fuerza de trabajo por sexo & DANE (GEIH) & \texttt{anex-GEIH}, hoja
Total nacional IML Sexo & Mensual \\
\texttt{empleo\_fuera\_ft.csv} & Población fuera de la fuerza de trabajo
por actividad principal & DANE (GEIH) & \texttt{anex-GEIH}, hoja
Población fuera de la fuerza de trabajo TN & Mensual \\
\texttt{empleo\_area.csv} & Indicadores de cabeceras y de centros
poblados y rural disperso & DANE (GEIH) & \texttt{anex-GEIH}, hoja Total
nacional Trim & Trimestre móvil \\
\texttt{empleo\_jovenes.csv} & Desempleo y participación de 15 a 28
años, proporción de jóvenes que no estudian ni están ocupados & DANE
(GEIH) & \texttt{anex-GEIHMLJ} & Trimestre móvil \\
\texttt{informalidad.csv} & Proporción de ocupados informales: nacional,
13 y 23 ciudades y áreas metropolitanas & DANE (GEIH) &
\texttt{anex-GEIHEISS} & Trimestre móvil (desde 2021) \\
\texttt{informalidad\_ramas.csv} & Ocupados, informales y tasa de
informalidad por rama & DANE (GEIH) & \texttt{anex-GEIHEISS}, hoja Ramas
de actividad & Trimestre móvil \\
\texttt{informalidad\_ciudades.csv} & Tasa de informalidad de 23
ciudades y áreas metropolitanas & DANE (GEIH) & \texttt{anex-GEIHEISS},
hoja Prop informalidad & Trimestre móvil \\
\texttt{inflacion\_clean.csv} & Inflación mensual y anual del IPC total
& BanRep (índice) y DANE (variaciones publicadas) & Serie 15000;
\texttt{anex-IPC} y \texttt{anex-IPC-Variacion} & Mensual \\
\texttt{ipc\_divisiones.csv} & Ponderación, variaciones y contribuciones
de las 12 divisiones & DANE & \texttt{anex-IPC}, hoja 2 & Último mes \\
\texttt{ipc\_ingresos.csv} & Variaciones por nivel de ingreso del hogar
& DANE & \texttt{anex-IPC}, hoja 3 & Último mes \\
\texttt{ipc\_ciudades.csv} & Variación anual por ciudad y división &
DANE & \texttt{anex-IPC}, hoja 6 & Último mes \\
\texttt{ipc\_subclases.csv} & Variaciones y contribución anual de las
188 subclases & DANE & \texttt{anex-IPC}, hoja 8 & Último mes \\
\texttt{ipc\_clasificaciones.csv} & Índices de energéticos, sin
alimentos ni energéticos, servicios, durables, semidurables y no
durables & DANE & \texttt{anex-IPC}, hojas 11 a 16 & Mensual (desde
2009) \\
\texttt{ipc\_ponderaciones.csv} & Ponderaciones oficiales de las 188
subclases por nivel de ingreso (canasta 2018) & DANE & Página ``IPC
ponderadores'' (actualización metodológica 2019) & Fija \\
\texttt{tasas\_interes\_clean.csv} & Curvas cero cupón de los TES en
pesos y en UVR a 1, 5 y 10 años & BanRep & 15272, 15273, 15274 (pesos);
15275, 15276, 15277 (UVR) & Diaria \\
\texttt{series\_banrep.csv} & Política monetaria, TRM, inflación básica,
externo y fiscal (detalle en 3.2) & BanRep, Superfinanciera, MinHacienda
y MinTrabajo vía BanRep & 16 series SUAMECA & Diaria a anual \\
\texttt{tasas\_mercado.csv} & IBR por plazos, TIB, DTF, CDT, tasas de
colocación, cartera y operaciones de liquidez (detalle en 3.2) & BanRep
(con información de la Superintendencia Financiera, formato 088) & 25
series SUAMECA & Diaria, semanal y mensual \\
\texttt{cambiario.csv} & Otras monedas, ITCR bilaterales y de
competitividad, balanza cambiaria, opciones PUT, dólar global, peso
mexicano y Brent (detalle en 3.2) & BanRep; Reserva Federal (H.10) y EIA
vía FRED & 18 series SUAMECA y 3 series FRED & Diaria, mensual y por
evento \\
\texttt{externo\_fiscal.csv} & Balanza de pagos, IED por sector,
remesas, deuda externa, PII, reservas, balance fiscal del GNC y del
SPNF, deuda del GNC, precios del comercio exterior (detalle en 3.2) &
BanRep; MinHacienda y DNP vía BanRep & 52 series SUAMECA & Mensual,
trimestral y anual \\
\texttt{colcap\_oficial.csv} & Índice COLCAP (puntos y base 100 en el 9
de febrero de 2009) & BVC/MSCI vía BanRep & Serie 6 & Diaria \\
\texttt{colcap\_canasta.csv} & Composición vigente del COLCAP: acciones,
emisor, sector y peso & \textbf{Auxiliar no oficial:} iShares MSCI
COLCAP (BlackRock) & Archivo de posiciones del fondo & Diaria
(última) \\
\texttt{acciones\_semanal.csv} & Cierres semanales de cada acción de la
canasta & \textbf{Auxiliar no oficial:} Yahoo Finance & Símbolos
\texttt{\textless{}TICKER\textgreater{}.CL} & Semanal (desde 2010) \\
\texttt{empresas\_10000\_agregados.csv} & Totales de las 10.000 empresas
más grandes por año, macrosector, región y supervisor, y concentración &
Supersociedades vía datos.gov.co & Conjunto \texttt{6cat-2gcs} & Anual
(cortes a diciembre) \\
\texttt{empresas\_10000\_top.csv} & Las 25 empresas con más ingresos de
cada año & Supersociedades vía datos.gov.co & Conjunto
\texttt{6cat-2gcs} & Anual \\
\texttt{empresas\_registro.csv} & Matrículas y cancelaciones mensuales
del registro mercantil por categoría (solo conteos agregados) &
Confecámaras/RUES vía datos.gov.co & Conjunto \texttt{c82u-588k} &
Mensual \\
\texttt{empresas\_banrep.csv} & Posición financiera neta por sector
institucional (\% del PIB) y deuda externa privada & BanRep & 16811,
16812, 16813, 16814 y 15332 & Trimestral y mensual \\
\texttt{exportaciones\_mensuales.csv} & Exportaciones FOB: café, carbón,
petróleo, ferroníquel, tradicionales, no tradicionales y total & DANE
con registros de la DIAN &
\texttt{anex-EXPORTACIONES-SerieCafeCarbonPetroleoNotradicionales} &
Mensual (desde 1992) \\
\texttt{exportaciones\_destinos.csv} & Exportaciones FOB por destino
principal & DANE con registros de la DIAN &
\texttt{anex-EXPORTACIONES-SeriePrincipalesDestinos} & Mensual \\
\texttt{importaciones\_mensuales.csv} & Importaciones CIF publicadas, de
zonas francas y total ampliado & DANE con registros de la DIAN &
\texttt{anex-IMP}, Cuadro A30 & Mensual \\
\texttt{importaciones\_cuode\_reciente.csv} & Importaciones por uso o
destino económico (CUODE): año corrido y último mes & DANE con registros
de la DIAN & \texttt{anex-IMP}, Cuadro A13 & Año corrido \\
\texttt{importaciones\_cuode\_anual.csv} & Importaciones CUODE por año &
DANE con registros de la DIAN & \texttt{anex-IMP-ImpoClasiCUODE} &
Anual \\
\texttt{importaciones\_origen.csv} & Importaciones CIF por país de
origen & DANE con registros de la DIAN &
\texttt{anex-IMP-MensPrincPaisesOrigen} & Mensual \\
\texttt{decisiones\_banrep.csv} & Última decisión de la Junta Directiva
ya anunciada y aún no vigente & Banco de la República (comunicado
oficial; registro manual) & --- & Por evento \\
\texttt{estado\_fuentes.csv} & Última observación y estado (vigente,
rezagado, pendiente) de cada fuente & Generado por el control de calidad
& --- & Cada corrida \\
\texttt{alertas\_datos.csv} & Observaciones TES aisladas marcadas como
sospechosas & Generado por el control de calidad & --- & Cada corrida \\
\end{longtable}\normalsize

Los anexos del DANE se descubren en la página oficial de cada operación
estadística; el proceso exige que exista exactamente un anexo vigente
con el patrón esperado y que su dominio sea el del DANE.

\hypertarget{series-del-banco-de-la-repuxfablica-por-archivo}{%
\subsubsection{3.2 Series del Banco de la República por
archivo}\label{series-del-banco-de-la-repuxfablica-por-archivo}}

\textbf{\texttt{series\_banrep.csv}.}

\small\begin{longtable}[]{@{}
  >{\raggedright\arraybackslash}p{(\columnwidth - 8\tabcolsep) * \real{0.2000}}
  >{\raggedright\arraybackslash}p{(\columnwidth - 8\tabcolsep) * \real{0.2000}}
  >{\raggedright\arraybackslash}p{(\columnwidth - 8\tabcolsep) * \real{0.2000}}
  >{\raggedright\arraybackslash}p{(\columnwidth - 8\tabcolsep) * \real{0.2000}}
  >{\raggedright\arraybackslash}p{(\columnwidth - 8\tabcolsep) * \real{0.2000}}@{}}
\toprule\noalign{}
\begin{minipage}[b]{\linewidth}\raggedright
Serie
\end{minipage} & \begin{minipage}[b]{\linewidth}\raggedright
Id SUAMECA
\end{minipage} & \begin{minipage}[b]{\linewidth}\raggedright
Entidad de origen
\end{minipage} & \begin{minipage}[b]{\linewidth}\raggedright
Frecuencia
\end{minipage} & \begin{minipage}[b]{\linewidth}\raggedright
Rango de control
\end{minipage} \\
\midrule\noalign{}
\endhead
\bottomrule\noalign{}
\endlastfoot
Tasa de política monetaria & 59 & BanRep & Diaria & 0 a 40 \\
IBR overnight, efectiva & 15324 & BanRep & Diaria & 0 a 40 \\
TRM & 1 & Superintendencia Financiera vía BanRep & Diaria & 500 a
10.000 \\
Inflación sin alimentos & 15388 & BanRep & Mensual & \(-5\) a 40 \\
Inflación sin alimentos ni regulados & 15390 & BanRep & Mensual & \(-5\)
a 40 \\
Inflación núcleo 15 & 15392 & BanRep & Mensual & \(-5\) a 40 \\
Inflación de regulados & 15398 & BanRep & Mensual & \(-20\) a 60 \\
Inflación de alimentos y bebidas & 15404 & DANE vía BanRep & Mensual &
\(-20\) a 60 \\
Meta de inflación & 853 & BanRep & Anual & 0 a 40 \\
ITCR-IPC, ponderaciones totales (2010 = 100) & 235 & BanRep & Mensual &
30 a 250 \\
Términos de intercambio & 15360 & BanRep & Mensual & 20 a 400 \\
Reservas internacionales netas & 15051 & BanRep & Mensual & \(-1.000\) a
500.000 \\
Cuenta corriente, \% del PIB & 15290 & BanRep & Trimestral & \(-20\) a
20 \\
Inversión extranjera directa en Colombia & 15133 & BanRep & Trimestral &
\(-20.000\) a 40.000 \\
Deuda bruta del GNC, \% del PIB & 15328 & MinHacienda vía BanRep & Anual
& 0 a 200 \\
Salario mínimo, variación anual & 15418 & MinTrabajo vía BanRep & Anual
& \(-5\) a 60 \\
\end{longtable}\normalsize

La meta de inflación, la inflación núcleo 15 y la variación del salario
mínimo se descargan y publican en el archivo, pero no alimentan cálculos
del sitio: el rango meta se fija en \(3\% \pm 1\) pp (sección 5.6).

\textbf{\texttt{tasas\_mercado.csv}.} IBR a 1, 3, 6 y 12 meses (15325,
15326, 16561, 16563); tasa interbancaria TIB (89); DTF a 90 días (65);
CDT a 90, 180 y 360 días (238, 239, 240); tasas de colocación total
(15110), sin tesorería (17280), consumo (15111), comercial ordinario
(15113), preferencial (15114), tesorería (15115), vivienda no VIS
(15105) y VIS (15107); cartera bruta en moneda legal total (373),
comercial (363), consumo (365), vivienda ajustada (371) y microcrédito
(367); saldos de repos de expansión a un día (17500), a otros plazos
(17501) y de la ventanilla de contracción (17503).

\textbf{\texttt{cambiario.csv}.} Del Banco de la República: BRL/USD
(15638), PEN/USD (15641), COP/EUR (30), COP/GBP (31), COP/JPY (33),
COP/CNY (28); ITCR bilateral con Estados Unidos (219), China (310),
Brasil (213) y México (226); ITCR-C (236); balanza cambiaria mensual:
cuenta corriente (16700), balanza comercial (16702), servicios y
transferencias (16704), movimientos netos de capital (16706), capital
del sector real y del Gobierno (16708) y variación de reservas brutas
(16712); monto aprobado en subastas de opciones PUT para acumulación de
reservas (16670). De FRED: índice amplio del dólar de la Reserva
Federal, H.10 (\texttt{DTWEXBGS}), pesos mexicanos por dólar, H.10
(\texttt{DEXMXUS}), y precio del Brent publicado por la EIA
(\texttt{DCOILBRENTEU}).

\textbf{\texttt{externo\_fiscal.csv}.} Cuenta corriente (15136), bienes
y servicios (15135), bienes (15706), servicios (15719), ingreso primario
(15140), ingreso secundario (15141), exportaciones (15707) e
importaciones (15708) de bienes; cuenta financiera total (16142),
inversión directa (16143), de cartera (16196), otra inversión (16303),
activos de reserva (16527), derivados (16271), errores y omisiones
(16544); IED en Colombia (15133) e IED de Colombia en el exterior
(15134); IED por sector: agro (15368), minería (15369), industria
(15370), electricidad (15371), construcción (15372), comercio (15373),
transporte (15374), financiero (15375), servicios (15376), petróleo
(15377); remesas (15363); deuda externa total (15330), pública (15331),
privada (15332) y como \% del PIB (15329); PII neta (16600), activos
(16601) y pasivos (16602); reservas (16850); Gobierno nacional central
en caja (mensual): ingresos (16722), gastos (16723), intereses (16724),
balance (16725), financiamiento interno (16726) y externo (16727);
sector público no financiero en caja (trimestral): 16728 a 16734; deuda
del GNC como \% del PIB (15328); índices de precios de exportación
(15361) e importación (15362) en dólares.

\hypertarget{identidad-y-vigencia-de-las-fuentes}{%
\subsubsection{3.3 Identidad y vigencia de las
fuentes}\label{identidad-y-vigencia-de-las-fuentes}}

Cada serie del graficador SUAMECA se descarga completa (toda la
historia, para incorporar revisiones). El cliente verifica que la
respuesta contenga una única serie, que su identificador sea el
solicitado y, cuando el código define un texto esperado, que el nombre
de la serie lo contenga (sin distinguir tildes ni mayúsculas). Si el
Banco reasigna un identificador, la descarga falla en lugar de publicar
otra variable. Las series de \texttt{externo\_fiscal.csv} y de
\texttt{empresas\_banrep.csv} se validan por identificador y por número
mínimo de observaciones, sin texto esperado (sección 9). La conexión al
servidor del Banco completa la cadena de certificados con el certificado
intermedio versionado en el repositorio, sin desactivar la verificación
TLS.

\hypertarget{fuentes-auxiliares-no-oficiales}{%
\subsubsection{3.4 Fuentes auxiliares no
oficiales}\label{fuentes-auxiliares-no-oficiales}}

La sección ``la bolsa por dentro'' (páginas Empresas y Portada) necesita
la composición del índice COLCAP y los precios de cada acción, que no se
obtienen en forma abierta y automatizable desde una fuente oficial. Se
usan dos fuentes auxiliares, rotuladas como tales en el sitio:

\small\begin{longtable}[]{@{}
  >{\raggedright\arraybackslash}p{(\columnwidth - 4\tabcolsep) * \real{0.3333}}
  >{\raggedright\arraybackslash}p{(\columnwidth - 4\tabcolsep) * \real{0.3333}}
  >{\raggedright\arraybackslash}p{(\columnwidth - 4\tabcolsep) * \real{0.3333}}@{}}
\toprule\noalign{}
\begin{minipage}[b]{\linewidth}\raggedright
Fuente
\end{minipage} & \begin{minipage}[b]{\linewidth}\raggedright
Uso
\end{minipage} & \begin{minipage}[b]{\linewidth}\raggedright
Limitación
\end{minipage} \\
\midrule\noalign{}
\endhead
\bottomrule\noalign{}
\endlastfoot
iShares MSCI COLCAP (BlackRock), archivo diario de posiciones del fondo
& Canasta y pesos de las acciones; selección de las ``7 Magníficas'';
peso por sector & Son los pesos del fondo (sin efectivo), una
aproximación a los pesos oficiales de la BVC/MSCI, no los pesos
oficiales \\
Yahoo Finance, cierres semanales de
\texttt{\textless{}TICKER\textgreater{}.CL} & Índices equiponderado y
``7 Magníficas''; variación de 12 meses de cada acción & Proveedor
comercial sin certificación oficial; precios sin ajuste por dividendos
ni eventos corporativos; posibles errores puntuales de precio \\
\end{longtable}\normalsize

El índice COLCAP mismo (nivel, variaciones, base 100) sí es la serie
oficial publicada por el Banco de la República (serie 6). Ninguna cifra
de las demás páginas depende de las fuentes auxiliares.

\hypertarget{registro-manual-de-decisiones-anunciadas}{%
\subsubsection{3.5 Registro manual de decisiones
anunciadas}\label{registro-manual-de-decisiones-anunciadas}}

La tasa de política del Banco de la República cambia en la serie oficial
el día en que empieza a regir, no el día del anuncio. Para no mostrar
una tasa ya superada durante ese intervalo,
\texttt{decisiones\_banrep.csv} registra manualmente la última decisión
anunciada (fecha del anuncio, fecha de vigencia, tasa y enlace al
comunicado). La regla de uso es estricta: la decisión se muestra como
``anunciada'' solo mientras la serie oficial no tenga datos en o después
de la fecha de vigencia; en cuanto los tenga, manda la serie oficial,
coincida o no con el registro manual. El registro no altera ningún
cálculo histórico (la tasa real, los ciclos y el traspaso usan la serie
oficial).

\hypertarget{transformaciones-comunes}{%
\subsection{4. Transformaciones
comunes}\label{transformaciones-comunes}}

Las transformaciones siguientes se usan en varias páginas; la sección 5
indica dónde.

\textbf{Variación anual y anualizaciones.}
\[\%\Delta_{a} x_t = 100\left(\frac{x_t}{x_{t-k}}-1\right),\qquad
\text{SAAR}_t = 100\left[\left(\frac{S_t}{S_{t-1}}\right)^{4}-1\right],\qquad
r^{3m}_t = 100\left[\left(\frac{\bar S_{t}}{\bar S_{t-3}}\right)^{4}-1\right],\]
donde \(S\) es una serie desestacionalizada y
\(\bar S_t = \tfrac13(S_t+S_{t-1}+S_{t-2})\). La inflación anualizada de
tres meses del IPC se calcula como
\(100\,[\prod_{j=0}^{2}(1+\pi^m_{t-j}/100)^{4}-1]\) sobre variaciones
mensuales sin desestacionalizar.

\textbf{Sumas y promedios móviles.} Las sumas de 4 trimestres o de 12
meses (\(X^{(4)}_t=\sum_{j=0}^{3}X_{t-j}\),
\(X^{(12)}_t=\sum_{j=0}^{11}X_{t-j}\)) eliminan la estacionalidad de
flujos sin necesidad de modelos y se usan para flujos externos,
fiscales, de comercio y sectoriales. Los promedios móviles de 3 o 12
meses suavizan tasas mensuales de la GEIH no desestacionalizadas. El
crecimiento ``de 12 meses'' compara la suma de los últimos cuatro
trimestres con la de los cuatro anteriores.

\textbf{Tasa de crecimiento compuesta.} Para un periodo de años
completos \(A_0,\dots,A_1\) con \(n=A_1-A_0+1\):
\(c = 100\,[(Y_{A_1}/Y_{A_0-1})^{1/n}-1]\), con \(Y_A\) la suma de los
cuatro trimestres del año.

\textbf{Índices base 100.} \(I_t = 100\,X_t/X_{b}\) con base fija (por
ejemplo, 2019-T4 = 100 o 2015 = 100). En los gráficos marcados como
``re-basados'', el navegador fija la base en el primer dato visible del
horizonte que elige el usuario (3, 5, 10 años o toda la historia).

\textbf{Deflactación.} Las magnitudes nominales se llevan a términos
reales con Fisher exacto: \(g^{r} = 100\,[(1+g)/(1+\pi/100)-1]\), donde
\(g\) es la variación nominal anual (en tanto por uno) y \(\pi\) la
inflación anual del IPC del mismo mes. Para ingresos anuales de empresas
se usa el promedio de las inflaciones anuales mensuales del año.

\textbf{Conversión a dólares.} El PIB nominal trimestral se convierte
con la TRM promedio del trimestre, \(Y^{\$}_q = Y_q/\overline{TRM}_q\),
y se suma en cuatro trimestres para las razones externas; el PIB anual
por persona en dólares usa la TRM promedio del año calendario.

\textbf{Participaciones y aportes.} Participación:
\(s_{i,t} = 100\,X_{i,t}/\sum_j X_{j,t}\). Aporte de un componente al
crecimiento anual de un agregado:
\(a_{i,t} = 100\,(X_{i,t}-X_{i,t-4})/X_{t-4}\), equivalente a
\(g_{i,t}\,w_{i,t-4}\). Con volúmenes encadenados los aportes no suman
exactamente el total; la diferencia se muestra (gasto) o se advierte
(oferta).

\textbf{Herfindahl y número equivalente.} \(H=\sum_i s_i^2\) con \(s_i\)
en tanto por uno y \(N^{eq}=1/H\): el número de categorías de igual
tamaño que daría la misma concentración.

\textbf{Volatilidad anualizada.} Desviación estándar móvil de 60
observaciones (mínimo 48) de los cambios diarios, multiplicada por
\(\sqrt{252}\): cambios logarítmicos en porcentaje para tasas de cambio;
cambios en puntos básicos para tasas de interés.

\textbf{Correlación móvil.} Correlación de Pearson de 52 semanas (mínimo
41) entre cambios porcentuales semanales (último dato de cada viernes).

\textbf{Percentil histórico.} Para cada indicador de la tabla de
señales, \(p = 100\cdot\#\{s\ge 2010: x_s \le x_T\}/\#\{s\ge 2010\}\),
con la muestra desde el 1 de enero de 2010.

\textbf{Recortes gráficos.} Algunos gráficos recortan la escala visible
(por ejemplo, \(\pm 12\%\) en el mapa de calor sectorial o \(-8\%\) en
la brecha mensual) para que 2020 no aplaste el resto de la historia; el
recuadro flotante muestra el valor exacto y ningún cálculo usa el valor
recortado.

\hypertarget{muxe9todos-por-puxe1gina}{%
\subsection{5. Métodos por página}\label{muxe9todos-por-puxe1gina}}

\hypertarget{portada}{%
\subsubsection{5.1 Portada}\label{portada}}

La portada no dibuja gráficos de Plotly: presenta la fase del ciclo
(sección 5.2), seis tarjetas con su estado y una minicurva (Crecimiento,
Inflación, Tasa del Banco, Desempleo, Dólar y Bolsa), una curva TES de
hoy frente a hace un año en SVG, un resumen por tema que toma la primera
frase de ``lo clave'' de cada página, la tabla de 16 indicadores y las
últimas publicaciones oficiales. El párrafo superior es un resumen en
lenguaje llano generado por reglas (sección 6). La forma de la curva se
clasifica con la pendiente \(S=y_{10}-y_{1}\): ``normal'' si \(S>0{,}3\)
pp, ``invertida'' si \(S<0\) y ``plana'' en otro caso.

\hypertarget{ciclo}{%
\subsubsection{5.2 Ciclo}\label{ciclo}}

\textbf{Brecha del producto principal: Hodrick-Prescott en tiempo real.}
Sea \(y_t=100\ln Y_t\) con \(Y_t\) el PIB real desestacionalizado del
DANE (Cuadro 4). La tendencia de Hodrick y Prescott (1997) resuelve
\[\min_{\tau}\ \sum_t (y_t-\tau_t)^2+\lambda\sum_t(\Delta^2\tau_t)^2,\qquad \lambda=1600,\]
cuya solución es el sistema lineal \((I+\lambda D'D)\,\tau = y\), con
\(D\) la matriz de segundas diferencias. La versión \textbf{en tiempo
real} (una cola) resuelve el problema, para cada trimestre \(t\), solo
con \(y_1,\dots,y_t\) y conserva el último punto \(\tau_{t\mid t}\); se
exige un mínimo de 20 trimestres. Así se evita el uso de información
futura y el sesgo de fin de muestra del filtro de dos colas (Orphanides
y van Norden, 2002). La brecha es \(g_t = y_t-\tau_{t\mid t}\).

\textbf{Tratamiento de la pandemia.} Los trimestres 2020-T2 a 2021-T2
(confinamiento a paro nacional) se reemplazan por interpolación lineal
\textbf{solo para estimar tendencias y coeficientes}; la brecha se mide
siempre contra el dato observado. Sin este tratamiento la tendencia
absorbería la caída de 2020 y deformaría la lectura de los años
siguientes.

\textbf{Cinco métodos y consenso.} Además de la brecha principal se
calculan:

\begin{itemize}
\tightlist
\item
  \textbf{HP de dos colas}: \(g^{2c}_t=y_t-\tau_{t\mid T}\), la mejor
  estimación ex post con toda la muestra.
\item
  \textbf{Hamilton (2018)}: residuo de la regresión por mínimos
  cuadrados ordinarios
  \[\tilde y_t = \beta_0+\sum_{k=0}^{3}\beta_{k+1}\,\tilde y_{t-8-k}+v_t,\qquad
  c^{H}_t = y_t-\Big(\hat\beta_0+\sum_{k=0}^{3}\hat\beta_{k+1}\,\tilde y_{t-8-k}\Big),\]
  con horizonte de 8 trimestres y 4 rezagos; \(\tilde y\) es la serie
  con la pandemia interpolada y la brecha se mide contra el \(y\)
  observado.
\item
  \textbf{Christiano y Fitzgerald (2003)}: filtro asimétrico de paso de
  banda para un paseo aleatorio con deriva, que aísla ciclos de 6 a 32
  trimestres (1,5 a 8 años, la banda de Baxter y King, 1999). Se retira
  la deriva lineal, se aplica el filtro sobre la serie interpolada y la
  tendencia es \(\tau^{CF}=\tilde y-CF_{6\text{–}32}(\tilde y)\); la
  brecha es \(y-\tau^{CF}\).
\item
  \textbf{Beveridge y Nelson (1981)}: con un AR(4) estimado por MCO para
  el crecimiento desviado de su media, \(x_t=\Delta\tilde y_t-\mu\), y
  su matriz compañera \(F\), el componente cíclico es
  \[c^{BN}_t = -\,e_1'\,F\,(I-F)^{-1}X_t,\qquad X_t=(x_t,\dots,x_{t-3})',\]
  evaluado sobre el crecimiento observado; es menos el crecimiento
  futuro esperado por encima de la media.
\end{itemize}

El consenso reporta la mediana, el mínimo, el máximo y el número de
métodos con brecha positiva. Solo la brecha HP en tiempo real está libre
de información futura; los otros cuatro métodos usan coeficientes o
filtros estimados con toda la muestra.

\textbf{Reloj del ciclo.} Siguiendo el reloj del ciclo de la OCDE, la
fase combina el nivel de la brecha y su dirección, medida como el cambio
en dos trimestres \(\Delta_2 g_t=g_t-g_{t-2}\) (reduce el ruido de un
solo dato):

\small\begin{longtable}[]{@{}lll@{}}
\toprule\noalign{}
& \(\Delta_2 g_t \ge 0\) & \(\Delta_2 g_t < 0\) \\
\midrule\noalign{}
\endhead
\bottomrule\noalign{}
\endlastfoot
\(g_t \ge 0\) & Expansión & Desaceleración \\
\(g_t < 0\) & Recuperación & Contracción \\
\end{longtable}\normalsize

La misma regla aplicada a la brecha de Hamilton produce una fase
alternativa con la que se mide la coincidencia entre métodos. La cinta
de fases colorea cada trimestre desde el primero con brecha en tiempo
real; la ``racha'' cuenta los trimestres consecutivos en la fase actual.

\textbf{Ciclo mensual.} Con el ISE desestacionalizado,
\(b_t = 100\ln ISE_t-\tau_t\), con \(\tau\) el HP en tiempo real de
frecuencia mensual con \(\lambda=129.600\) (equivalente mensual de 1.600
según Ravn y Uhlig, 2002), mínimo 36 meses y los meses de marzo de 2020
a junio de 2021 interpolados para estimar la tendencia. La fase mensual
usa \(\Delta_3 b_t=b_t-b_{t-3}\). El mismo cálculo se aplica a las tres
grandes ramas del ISE (primarias, secundarias y terciarias), que se
muestran junto a su variación anual.

\textbf{Motores sectoriales.} Los aportes de los 12 sectores al
crecimiento anual del valor agregado son los de la sección 5.3. La
amplitud cuenta cuántos de los 12 sectores producen por encima de su
propia tendencia (brecha sectorial de la sección 5.4).

\textbf{Ritmo.} Crecimiento anual del ISE desestacionalizado frente al
ritmo de los últimos tres meses anualizado (\(r^{3m}\) de la sección 4).
Si el segundo supera al primero, el texto dice que la actividad acelera.

\textbf{Puntos de giro del ISE.} En la línea de Bry y Boschan (1971),
sobre el promedio móvil centrado de tres meses
\(S_t=\tfrac13(ISE_{t-1}+ISE_t+ISE_{t+1})\) se marca un pico (valle)
cuando \(S_t\) es el máximo (mínimo) de la ventana \([t-6,\,t+6]\); los
giros se fuerzan a alternar (entre dos picos seguidos se conserva el más
alto y entre dos valles el más bajo) y se descarta un valle inicial
seguido por un pico a menos de seis meses. Una recesión va de pico a
valle y una expansión de valle a pico (o hasta el último dato si sigue
en curso); para cada una se reportan duración en meses y variación de
\(S\). Las recesiones sombrean varios gráficos del sitio. Este fechado
es descriptivo: el promedio centrado y la ventana usan meses posteriores
a cada giro, por lo que los giros más recientes pueden cambiar con datos
nuevos.

\textbf{Empleo y ley de Okun.} El cambio en 12 meses del promedio de
tres meses de las tasas de desempleo y ocupación desestacionalizadas
resume la respuesta laboral. Para Okun (1962) en brechas, el desempleo
desestacionalizado se promedia por trimestre (\(u_q\)), se estima su
tendencia \(u^*_q\) con HP de dos colas (\(\lambda=1600\), pandemia
interpolada) y se define \(b^u_q=u_q-u^*_q\); luego
\[b^u_q = \alpha+\beta\, g^{2c}_q+\varepsilon_q\] por MCO excluyendo
2020 y 2021, con \(g^{2c}\) la brecha HP de dos colas. Se reportan
\(\hat\beta\), la correlación y el número de trimestres. El gráfico
muestra \(-b^u\) para que ambas series suban juntas.

\hypertarget{crecimiento-incluye-sectores-gasto-inversiuxf3n-consumo-y-pib-por-persona}{%
\subsubsection{5.3 Crecimiento (incluye sectores, gasto, inversión,
consumo y PIB por
persona)}\label{crecimiento-incluye-sectores-gasto-inversiuxf3n-consumo-y-pib-por-persona}}

\textbf{Tres velocidades.} (i) Variación anual del PIB real original,
\(100(Y_t/Y_{t-4}-1)\); (ii) crecimiento de 12 meses, suma de cuatro
trimestres frente a los cuatro anteriores; (iii) variación trimestral
anualizada del PIB desestacionalizado (SAAR). También se informa la
variación trimestral sin anualizar.

\textbf{Año corrido y crecimiento anual.} El año corrido suma los
trimestres publicados del año en curso y los compara con los mismos
trimestres del año anterior. El crecimiento por año calendario compara
la suma de los cuatro trimestres de cada año con la del año anterior; si
el año en curso está incompleto, se muestra la barra de los últimos 12
meses. Se acompaña del ``ritmo habitual'' (promedio anual del
crecimiento de la tendencia HP de dos colas,
\(g^*_t=100\,[e^{(\tau_t-\tau_{t-4})/100}-1]\)) y del promedio simple
2010--2019. La tendencia de dos colas se usa para el crecimiento
potencial porque la de tiempo real oscila con el rebote de 2021--2022.

\textbf{Crecimiento por periodos.} Tasas compuestas para ``Auge
petrolero'' (2010--2014), ``Ajuste'' (2015--2019), ``Pandemia y rebote''
(2020--2021) y ``Pospandemia'' (2022--2025), más el año corrido. Un
periodo se incluye solo si sus años están completos.

\textbf{Real frente a nominal y deflactor.} El deflactor implícito es
\(D_t = N_t/Y_t\) (PIB nominal sobre real, datos originales) y su
inflación es \(\pi^D_t=100(D_t/D_{t-4}-1)\), de modo que
\((1+n)=(1+g)(1+\pi^D)\). Se compara con el IPC promediado por
trimestre. La divergencia entre deflactor e IPC se interpreta, siguiendo
a Kohli (2004), como efecto de los términos de intercambio; el texto
reporta su variación anual (promedio trimestral de la serie 15360).

\textbf{Sectores (oferta).} Para las 12 agrupaciones CIIU Rev.~4
(volúmenes encadenados, datos originales): variación anual \(g_{i,t}\),
peso \(w_{i,t}=100\,X_{i,t}/VA_t\) y aporte
\(a_{i,t}=100\,(X_{i,t}-X_{i,t-4})/VA_{t-4}\), con \(VA\) el valor
agregado total. Los aportes se agrupan en primario (agro y minería),
industria-energía-construcción, servicios de mercado (comercio,
transporte y turismo; comunicaciones; finanzas; inmobiliarias; servicios
profesionales; arte y hogares) y Gobierno, educación y salud. El
crecimiento ``sin el sector \(X\)'' (Gobierno o minería) es
\[g_{-X,t}=100\,\frac{\sum_i a_{i,t}-a_{X,t}}{\sum_i w_{i,t}-w_{X,t}},\]
una aproximación que usa los pesos corrientes del trimestre. El índice
de difusión cuenta los sectores con \(g_{i,t}>0\); la comparación con la
historia contrasta el promedio de \(g_{i,t}\) en 2015--2019 con el de
los últimos cuatro trimestres. El mapa de calor recorta la escala de
color en \(\pm 12\%\).

\textbf{Nivel frente al camino previo.} Se ajusta por MCO
\(\ln S_k = b_0+b_1 k\) sobre el PIB desestacionalizado de 2015-T1 a
2019-T4 y se extiende \(\hat S_k=e^{b_0+b_1k}\). El ritmo anual del
camino es \(100\,(e^{4b_1}-1)\) y la distancia es
\(100\,(S_t/\hat S_t-1)\). El camino es una referencia contrafactual
descriptiva (Cerra y Saxena, 2008; Blanchard, Cerutti y Summers, 2015),
no una proyección.

\textbf{Demanda (gasto).} Con los componentes en volúmenes encadenados,
datos originales: aportes de consumo de los hogares, Gobierno e
inversión fija, \(a_{k,t}=100\,(C_{k,t}-C_{k,t-4})/PIB_{t-4}\); comercio
exterior neto, \(a_X-a_M\); y un residuo ``existencias y discrepancia''
igual al crecimiento del PIB menos la suma anterior. También se reporta
la variación anual de cada componente y de la demanda interna. El
proceso verifica que el PIB del anexo de gasto coincida con el del anexo
de producción (diferencia relativa máxima de 0,2\%).

\textbf{Inversión.} Tasa de inversión
\(TI_t=100\,\sum_{j=0}^{3}FBKF^{nom}_{t-j}/\sum_{j=0}^{3}PIB^{nom}_{t-j}\)
y, por tipo de activo (vivienda, otros edificios y obras, maquinaria y
equipo, propiedad intelectual), índice de la inversión real
desestacionalizada con 2019-T4 = 100.

\textbf{Consumo de los hogares.} Variación de 12 meses (suma de cuatro
trimestres frente a los cuatro previos) por durabilidad (durables,
semidurables, no durables, servicios) y por finalidad COICOP (12
grupos); la participación aproximada de cada finalidad usa la suma de
volúmenes encadenados de los últimos cuatro trimestres. La participación
del consumo en el PIB usa valores corrientes de 12 meses.

\textbf{PIB por persona y productividad.} PIB real anual por persona
\(y_A=\sum_{q\in A}Y_q/P_A\) (pesos de 2015) con \(P_A\) la población a
mitad de año, solo para años completos; en dólares,
\(\sum_{q\in A}N_q/(P_A\,\overline{TRM}_A)\), que mezcla crecimiento,
inflación y tipo de cambio. La productividad laboral aparente (OECD,
2001) es el cociente de dos índices 2015 = 100: PIB real de cuatro
trimestres y promedio de 12 meses de los ocupados desestacionalizados;
su variación anual compara el último trimestre con el mismo del año
anterior.

\hypertarget{capacidad-incluye-regiones-y-departamentos}{%
\subsubsection{5.4 Capacidad (incluye regiones y
departamentos)}\label{capacidad-incluye-regiones-y-departamentos}}

\textbf{Producción frente a capacidad.} La capacidad (PIB potencial) se
obtiene de la brecha HP en tiempo real: \(Y^*_t=Y_t\,e^{-g_t/100}\). El
gráfico muestra ambos niveles en billones de pesos de 2015 por trimestre
y colorea la distancia (verde por encima, naranja por debajo). Las
tarjetas resumen la mediana y el rango de los cinco métodos de la
sección 5.2.

\textbf{Brecha por sectores.} Para cada uno de los 12 sectores, el nivel
se anualiza con la suma de cuatro trimestres (elimina la estacionalidad
de los datos originales sin modelos),
\(N_{i,t}=\sum_{k=0}^{3}X_{i,t-k}\), y se calcula
\(b_{i,t}=100\ln N_{i,t}-\tau_{i,t}\) con HP en tiempo real
(\(\lambda=1600\), pandemia interpolada). El mapa de calor desde 2016
ordena los sectores por su brecha más reciente.

\textbf{Holgura laboral.} Tres medidas de la OIT (2013), promediadas en
12 meses: tasa de desempleo \(TD=D/FT\); desempleo más subocupación por
horas \(TCSD=(D+S)/FT\); y la medida compuesta
\(MCSFT=(D+S+FTP)/(FT+FTP)\), con \(FTP\) la fuerza de trabajo
potencial. La brecha de desempleo frente a su tendencia es la de la ley
de Okun (sección 5.2). En \texttt{laboral\_subutilizacion.csv} los ceros
del anexo se tratan como dato faltante.

\textbf{Regiones.} Con el PIB departamental del DANE (base 2015):
crecimiento real del último año publicado, variación frente a 2019,
\(100\,(Y_{d,A}/Y_{d,2019}-1)\), y PIB por habitante relativo,
\(100\,pc_d/pc_{COL}\) (pesos corrientes). El mapa esquemático asigna
cada medida a siete tramos de color con cortes fijos. El desempleo de la
ciudad capital (año móvil de la GEIH) completa el mapa. El proceso
verifica que la suma de los departamentos reproduzca el total nacional
con una diferencia menor al 1\%.

\textbf{Estructura y diversificación.} Participación de cada una de las
12 ramas en el valor agregado corriente del departamento (sin impuestos)
y número equivalente de sectores \(1/\sum_i s_{i,d}^2\) (Herfindahl,
1950; Hirschman, 1964), que va de 1 a 12.

\textbf{Desempleo en las 32 capitales.} Tasa del año móvil más reciente
frente al año móvil terminado un año antes (se presenta en la página
Empleo).

\hypertarget{empleo}{%
\subsubsection{5.5 Empleo}\label{empleo}}

\textbf{Desempleo e informalidad.} La tasa de desempleo
desestacionalizada del DANE se muestra mensual y en promedio de tres
meses, con su cambio frente a un año antes. La informalidad es la
proporción de ocupados informales de la GEIH (definición del DANE
alineada con la OIT), en trimestres móviles desde 2021, a nivel
nacional, en 13 y 23 ciudades, por rama (\(I_r/O_r\)) y por ciudad. El
lector verifica que la secuencia de trimestres móviles sea consecutiva
(los rótulos del DANE no siempre traen el año). La serie no es
comparable con la medición anterior a 2021.

\textbf{Series sin desestacionalizar.} Los anexos de detalle de la GEIH
no están desestacionalizados; por eso las comparaciones usan promedios
de 12 meses o el mismo periodo del año anterior:

\begin{itemize}
\tightlist
\item
  Empleos creados o perdidos por rama y por posición:
  \(\bar O_t-\bar O_{t-12}\), con \(\bar O\) el promedio de tres meses,
  en miles de personas.
\item
  Empleo por rama frente a 2019: promedio de los últimos 12 meses frente
  al promedio de 2019.
\item
  Composición por posición ocupacional (CISE-93, OIT 1993):
  participación sobre el promedio de 12 meses de las siete posiciones
  graficadas; los asalariados suman las posiciones privada y del
  Gobierno.
\item
  Brechas por sexo: tasas de desempleo y de participación promediadas en
  12 meses; brecha mujeres − hombres en pp.
\item
  Jóvenes de 15 a 28 años (Ley 1622 de 2013): desempleo por sexo y
  proporción de jóvenes que no estudian ni están ocupados, en trimestre
  móvil; la serie de mujeres y la de hombres suman el total.
\item
  Cabeceras frente a centros poblados y rural disperso: desempleo en
  trimestre móvil.
\item
  Población fuera de la fuerza de trabajo por actividad principal:
  promedio de 12 meses, en millones.
\end{itemize}

\textbf{Ocupados totales.} La tarjeta de ocupados suma las ramas y
promedia 12 meses; su cambio compara con el mismo promedio un año antes.

\hypertarget{inflaciuxf3n}{%
\subsubsection{5.6 Inflación}\label{inflaciuxf3n}}

\textbf{Inflación total.} El índice del IPC (serie 15000 del Banco de la
República) se descarga completo y se valida que no tenga meses
faltantes. La inflación anual es \(\pi_t=100\,(IPC_t/IPC_{t-12}-1)\) y
la mensual \(100\,(IPC_t/IPC_{t-1}-1)\), redondeadas a dos decimales.
Las variaciones mensuales se reemplazan por las publicadas por el DANE
en su tabla histórica cuando existen (se exige que no difieran más de
0,10 pp del cálculo con el índice) y la variación anual y mensual del
último mes se reemplaza por la cifra oficial del DANE (diferencia máxima
tolerada de 0,05 pp; el índice publicado por el Banco tiene dos
decimales). El rango meta del Banco es \(3\% \pm 1\) pp.

\textbf{Inflación de fondo y por grupos.} Inflación sin alimentos ni
regulados (serie 15390, la medida básica principal), sin alimentos
(15388), de regulados (15398) y de alimentos (15404), en nivel y frente
a un año antes.

\textbf{Dirección y momento.} El cambio de la inflación anual en tres
meses, \(\pi_t-\pi_{t-3}\), define si acelera o cede (sección 6). La
inflación anualizada de tres meses se calcula sin desestacionalizar y se
advierte la estacionalidad de enero a marzo.

\textbf{Inflación implícita (breakeven).} Con las curvas cero cupón de
los TES en pesos (\(i_h\)) y en UVR (\(r_h\)):
\[\pi^{BE}_h = 100\left[\frac{1+i_h/100}{1+r_h/100}-1\right],\qquad h\in\{1,5,10\}.\]
El rezago de indexación de la UVR (cerca de un mes) se ignora para
plazos de un año o más.

\textbf{Forward 5y5y.} Con factores forward nominal y real,
\[F^{\$}=\left[\frac{(1+i_{10})^{10}}{(1+i_5)^{5}}\right]^{1/5},\qquad
F^{UVR}=\left[\frac{(1+r_{10})^{10}}{(1+r_5)^{5}}\right]^{1/5},\qquad
\pi^{5y5y}=100\,\big(F^{\$}/F^{UVR}-1\big),\] equivalente a
\([(1+b_{10})^{10}/(1+b_5)^5]^{1/5}-1\) en términos de los breakevens
\(b_h\).

\textbf{Trayectoria que descuenta el mercado.} Con \(b_1,b_5,b_{10}\):
año 1, \(f_{0,1}=b_1\); años 1 a 5,
\(f_{1,5}=[(1+b_5)^5/(1+b_1)]^{1/4}-1\); años 5 a 10,
\(f_{5,10}=[(1+b_{10})^{10}/(1+b_5)^5]^{1/5}-1\). Los tres tramos
encadenan exactamente el breakeven a 10 años,
\((1+f_{0,1})(1+f_{1,5})^4(1+f_{5,10})^5=(1+b_{10})^{10}\). El gráfico
los dibuja hacia adelante desde la última fecha, junto a la misma
lectura de hace un año. Es un precio observado hoy, no un pronóstico del
tablero.

\textbf{Aportes por división y subclase.} Las contribuciones a la
inflación anual son las que publica el DANE,
\(A_{i,t}=w_i\,(I_{i,t-12}/I_{t-12})\,\pi_{i,t}\), con \(w_i\) la
ponderación de la canasta 2018; suman la inflación total.

\textbf{Bienes y servicios.} Variación anual de los índices del DANE de
servicios, no durables, semidurables, durables, energéticos y sin
alimentos ni energéticos (Bryan y Cecchetti, 1994; Baumol, 1967).

\textbf{Difusión.} Proporción sin ponderar de las 188 subclases con
inflación anual mayor a 4\% (techo del rango) y mayor a 6\%, y la misma
proporción ponderada por el peso de cada subclase en la canasta
(Cecchetti, 1997). Las subclases se unen a sus ponderaciones por
posición cuando las dos listas coinciden en al menos 95\% de los
nombres, y por nombre en caso contrario. El histograma usa intervalos de
1 pp con los extremos agrupados en \(-10\%\) y \(20\%\).

\textbf{Ingresos y ciudades.} IPC calculado por el DANE con la canasta
de cada nivel de ingreso (pobres, vulnerables, clase media, ingresos
altos; Jaravel, 2021) y por ciudad (23 ciudades más otras áreas urbanas)
y división.

\hypertarget{tasas-de-interuxe9s}{%
\subsubsection{5.7 Tasas de interés}\label{tasas-de-interuxe9s}}

\textbf{Tasa real ex ante y ex post.} Con la tasa de política
\(i^{TPM}\) y el breakeven a un año:
\[r^{ea}_t = 100\left[\frac{1+i^{TPM}_t/100}{1+\pi^{BE}_{1,t}/100}-1\right],\qquad
r^{ep}_t = 100\left[\frac{1+i^{TPM}_t/100}{1+\pi^{pub}_t/100}-1\right],\]
donde \(\pi^{pub}_t\) es la inflación anual del último mes
\textbf{publicado} a la fecha \(t\): el IPC del mes \(m\) se considera
disponible desde el día 11 del mes \(m+1\) (un mes más diez días, regla
conservadora), de modo que no se usa información futura. La tasa de
política vigente cada día se asigna con el último dato disponible (unión
hacia atrás).

\textbf{Tasa neutral.} Se usa como parámetro externo el rango de 2,7\% a
3,0\% real, estimación del equipo técnico del Banco de la República
(2025, con señal de ajuste al alza hacia 2026); conceptualmente, la tasa
que ni frena ni estimula la economía (Laubach y Williams, 2003). La
clasificación de la postura está en la sección 6.

\textbf{Ciclos de la tasa de política.} Una decisión es un día con
\(\Delta TPM\neq 0\). Un ciclo agrupa decisiones consecutivas del mismo
signo y termina cuando la siguiente decisión va en sentido contrario
(las pausas no lo cortan). Para cada ciclo desde 2000 se reportan
fechas, número de decisiones, tasa antes y después, cambio acumulado y
duración en meses, \(\text{round}(\text{días}/30{,}44)+1\).

\textbf{Traspaso.} Para cada uno de los seis ciclos más recientes desde
junio de 2008 y cada tasa de mercado \(k\) (IBR a un día y a 3 meses,
CDT y DTF a 90 días, colocación total, ordinaria, preferencial, consumo
y vivienda):
\[P_k = 100\,\frac{R_{k,t_1}-R_{k,t_0}}{TPM_{t_1}-TPM_{t_0}},\] con
\(t_0\) el día anterior a la primera decisión y \(t_1\) tres meses
después de la última (o el último dato). Si la serie no existía en
\(t_0\) se muestra ``---''. El texto compara los cambios en pp del
último ciclo completo y del ciclo actual.

\textbf{Costo del crédito y margen.} Último dato semanal de cada
modalidad de colocación y su tasa real con la inflación anual más
reciente (Fisher exacto). Margen de intermediación: promedio mensual de
la colocación total menos el del CDT a 90 días; prima del consumo:
consumo menos tasa de política.

\textbf{Mercado monetario.} Curva del IBR (un día, 1, 3, 6 y 12 meses)
hoy, hace tres meses y hace un año, con la tasa de política de cada
fecha. Diferenciales diarios del IBR a un día y de la TIB frente a la
tasa de política, en pb, promediados por semana; el texto reporta la
desviación absoluta media del último año.

\textbf{Cartera y liquidez.} Crecimiento real anual del saldo de cartera
en moneda legal por modalidad (deflactado con la inflación anual del
mes); composición sobre la suma de comercial, consumo, vivienda y
microcrédito. Liquidez: saldos diarios de repos de expansión (a un día y
a plazo) y de contracción, promediados por mes, en billones (los días
sin operación cuentan como cero).

\hypertarget{curva-tes}{%
\subsubsection{5.8 Curva TES}\label{curva-tes}}

\textbf{Datos.} Tasas cero cupón a 1, 5 y 10 años en pesos y en UVR que
estima el Banco de la República con el modelo de Nelson y Siegel (1987)
sobre operaciones del SEN y del MEC (Arango, Melo y Vásquez, 2002). El
explorador permite ver la curva de cualquier día desde 2003 y compararla
con cierres de años anteriores.

\textbf{Factores.} Nivel \(N_t=(y_1+y_5+y_{10})/3\); pendiente
\(S_t=y_{10}-y_1\); curvatura \(C_t=2y_5-y_1-y_{10}\) (Litterman y
Scheinkman, 1991). El percentil de la pendiente es la proporción de días
desde 2003 con pendiente mayor que la actual.

\textbf{Episodios de curva invertida.} Tramos de al menos cinco días
hábiles consecutivos con \(S_t<0\); para cada uno se reportan inicio,
fin, duración, pendiente mínima y tasa de política al inicio (Estrella y
Hardouvelis, 1991).

\textbf{Movimientos.} Cambio de cada plazo en 1, 3 y 12 meses (pb). El
movimiento de 12 meses se clasifica por el signo del cambio de nivel
(alza o baja) y de pendiente (empinamiento o aplanamiento); si ambos
cambios son menores a 0,10 pp en valor absoluto se rotula ``sin cambio
de fondo''. El cambio nominal de cada plazo se descompone en el cambio
de la tasa UVR y el cambio de la diferencia nominal menos real.

\textbf{Fisher, prima y volatilidad.} La compensación por inflación
exacta es \((1+y^{\$}_n)/(1+r_n)-1\) y se muestra junto a la diferencia
simple \(y^{\$}_n-r_n\). La prima sobre la tasa de política es
\(y_n-TPM\) para \(n=1\) y 10, con su promedio desde 2003. La
volatilidad anualizada del TES a 10 años es la desviación estándar móvil
de 60 días de los cambios diarios en pb por \(\sqrt{252}\).

\hypertarget{mercados-tasa-de-cambio-y-otras-monedas}{%
\subsubsection{5.9 Mercados: tasa de cambio y otras
monedas}\label{mercados-tasa-de-cambio-y-otras-monedas}}

\textbf{TRM.} Pesos por dólar certificados por la Superintendencia
Financiera; los gráficos usan el último dato de cada semana y su
variación frente a un año antes. Variaciones de 12 meses: último dato
frente al último disponible 12 meses antes.

\textbf{Peso frente al dólar global y sus pares.} Se comparan las
variaciones de 12 meses del peso colombiano, el real brasileño, el peso
mexicano y el sol peruano (unidades por dólar; positivo = la moneda se
debilita) con la del índice amplio del dólar de la Reserva Federal. La
diferencia frente a los pares es
\(d=v_{COP}-\tfrac13(v_{BRL}+v_{MXN}+v_{PEN})\). Se calcula también la
correlación de las variaciones anuales semanales de la TRM y del dólar
global.

\textbf{Otras monedas.} Tasas cruzadas
\(COP/X = TRM/(X\text{ por USD})\) para real, peso mexicano y sol; euro,
libra, yen y yuan directamente del Banco. Índices re-basados al
horizonte elegido y variación de 12 meses (negativo = el peso se
fortalece).

\textbf{Tasa de cambio real.} ITCR-IPC con ponderaciones de comercio
total e ITCR-C (competitividad en el mercado de EE. UU.), base 2010 =
100, y su promedio desde 2000; aumentos = depreciación real. Las cuatro
series bilaterales (Estados Unidos, China, Brasil, México) se expresan
como distancia porcentual frente a su propio promedio desde 2000; el
sitio las describe como deflactadas con el índice de precios del
productor. El Banco advierte que 2010 es una fecha de comparación, no un
nivel de equilibrio.

\textbf{Petróleo y términos de intercambio.} Con promedios mensuales, se
calculan las variaciones anuales del Brent (\(x\)) y de la TRM (\(y\)) y
se estiman por MCO, por separado para 2008--2019 y desde 2020, la recta
\(\hat y=a+bx\) y la correlación; el texto traduce \(b\) al efecto de un
alza de 10\% del Brent. La participación del petróleo y del carbón en
las exportaciones usa sumas de 12 meses en dólares FOB (Chen y Rogoff,
2003). Correlaciones móviles de 52 semanas de la TRM con el Brent y con
el dólar global.

\textbf{Flujos y reservas.} Balanza cambiaria mensual (cuenta corriente,
movimientos de capital y variación de reservas brutas) en sumas de 12
meses; reservas internacionales netas a fin de mes y monto aprobado en
subastas de opciones PUT para acumulación, sumado por año.

\textbf{Volatilidad.} Desviación estándar móvil de 60 observaciones de
los cambios logarítmicos diarios, anualizada, para el peso colombiano,
el real y el peso mexicano.

\hypertarget{empresas-bolsa-10.000-empresas-y-registro-mercantil}{%
\subsubsection{5.10 Empresas: bolsa, 10.000 empresas y registro
mercantil}\label{empresas-bolsa-10.000-empresas-y-registro-mercantil}}

\textbf{COLCAP oficial.} Índice de precios (sin dividendos) de la
BVC/MSCI publicado por el Banco de la República. La transición del
COLCAP de la BVC al MSCI COLCAP el 28 de mayo de 2021 mantuvo la
continuidad de niveles. Se calculan la variación de 12 meses (último
dato frente al último disponible 365 días antes), la caída desde el
máximo histórico (\(100\,(P_t/\max_{s\le t}P_s-1)\)) y el COLCAP en
dólares, puntos divididos por la TRM del mismo día o del último día
disponible con tolerancia de siete días.

\textbf{La bolsa por dentro (fuentes auxiliares, sección 3.4).}

\begin{itemize}
\tightlist
\item
  \emph{Canasta y emisores.} La canasta vigente del fondo iShares MSCI
  COLCAP (solo acciones colombianas con peso positivo) se agrupa por
  emisor uniendo acciones ordinarias y preferenciales; se exige que los
  pesos sumen entre 80\% y 101\%.
\item
  \emph{7 Magníficas.} Las siete empresas de mayor peso agregado en la
  canasta vigente; para cada una se usa la clase de acción de mayor
  peso. Se recalculan solas con cada cambio de canasta.
\item
  \emph{Precios.} Cierres semanales de Yahoo Finance fechados al
  viernes. Un ticker que cambió (por ejemplo, PFDAVVNDA a PFDAVIGRP) se
  empalma hacia atrás escalando la serie anterior en la primera semana
  común.
\item
  \emph{Índices.} Equiponderado con rebalanceo semanal,
  \[E_t = E_{t-1}\Big(1+\tfrac{1}{N_t}\textstyle\sum_{i} r_{i,t}\Big),\qquad r_{i,t}=P_{i,t}/P_{i,t-1}-1,\]
  donde se descartan los retornos con \(\lvert r_{i,t}\rvert>0{,}60\)
  como errores de la fuente y \(N_t\) es el número de retornos válidos;
  el índice de las 7 Magníficas aplica la misma regla a sus siete
  acciones. Las tres series (COLCAP semanal, equiponderado y 7
  Magníficas) se llevan a base 100 en la primera semana común y el
  navegador las re-basa al horizonte elegido.
\item
  \emph{Lectura.} Si en 12 meses el COLCAP supera al equiponderado por
  más de 2 pp, el avance se concentra en las empresas grandes; si el
  equiponderado supera al COLCAP por más de 2 pp, el avance es amplio.
\item
  \emph{Acciones y sectores.} Variación de 52 semanas de la acción
  principal de cada emisor (último cierre frente al último disponible
  364 días antes) y suma de pesos por sector.
\end{itemize}

\textbf{Las 10.000 empresas más grandes.} Del conjunto
\texttt{6cat-2gcs} de Supersociedades (cifras en billones de pesos
corrientes, cortes a diciembre) se agregan por año, macrosector, región
y supervisor los ingresos operacionales, la ganancia (pérdida), los
activos, los pasivos y el patrimonio; se exige que cada año tenga al
menos 9.000 empresas. Razones sobre sumas agregadas: margen neto
\(\sum G/\sum I\), rentabilidad del patrimonio \(\sum G/\sum PAT\) y
endeudamiento \(\sum P/\sum A\). Concentración: participación de las
\(N\) mayores en los ingresos del año para
\(N\in\{10,50,100,500,1.000,2.500,5.000,10.000\}\) (Gabaix, 2011). El
crecimiento real de los ingresos se deflacta con el promedio de las
inflaciones anuales mensuales del año. La lista cambia cada año, de modo
que las comparaciones son entre las 10.000 más grandes de cada año, no
entre las mismas empresas.

\textbf{Financiación de las empresas.} Crecimiento real de la cartera
comercial y total, tasas preferencial y ordinaria (promedios mensuales),
posición financiera neta por sector institucional (\% del PIB, cuentas
financieras del Banco), deuda externa privada e IED de cuatro
trimestres.

\textbf{Registro mercantil.} Del conjunto \texttt{c82u-588k} se
consultan \textbf{solo conteos agregados} por fecha y categoría de
matrícula, nunca registros individuales ni datos personales. Las
matrículas se agregan por mes de matrícula y las cancelaciones por mes
de cancelación, desde 2010. ``Sociedades'' agrupa las categorías que
empiezan por ``SOCIEDAD''; el resto se reporta como personas naturales y
otras categorías. Se descarta el último mes si sus matrículas son
menores a la mitad de la mediana de los 12 meses anteriores (mes
incompleto). Se reportan sumas de 12 meses y la variación de las
matrículas de sociedades frente a los 12 meses previos.

\hypertarget{externo-balanza-de-pagos-y-cuentas-fiscales}{%
\subsubsection{5.11 Externo: balanza de pagos y cuentas
fiscales}\label{externo-balanza-de-pagos-y-cuentas-fiscales}}

\textbf{Balanza de pagos (MBP6).} Los componentes de la cuenta corriente
(bienes, servicios, ingreso primario e ingreso secundario) se suman en
cuatro trimestres; su suma es la cuenta corriente. La cuenta financiera
usa el signo del MBP6 (activos netos menos pasivos netos; negativo =
entrada neta de capital); para graficar las entradas netas por tipo se
invierte el signo, \(F_i=-CF_i\), y se compara con el déficit corriente
a financiar \(-CC\). Exportaciones e importaciones de bienes de la
balanza de pagos (FOB) en cuatro trimestres.

\textbf{Razones al PIB en dólares.} \(CC^{(4)}/Y^{\$(4)}\), con
\(Y^{\$(4)}\) la suma de cuatro trimestres del PIB nominal del DANE
convertido con la TRM promedio de cada trimestre. Se reportan también la
cuenta corriente del último trimestre sobre el PIB en dólares de ese
trimestre, la balanza de bienes, la IED y la PII neta. La tarjeta
histórica de cuenta corriente de la portada usa la razón publicada por
el Banco (serie 15290).

\textbf{IED y remesas.} IED de cuatro trimestres y su cobertura del
déficit corriente, \(100\,IED^{(4)}/\lvert CC^{(4)}\rvert\) (si hay
déficit); IED por sector, últimos cuatro trimestres frente a los cuatro
anteriores. Remesas en sumas de 12 meses y como porcentaje del PIB en
dólares de los últimos cuatro trimestres disponibles.

\textbf{Deuda externa, PII y reservas.} Saldos de deuda externa pública
y privada (la razón al PIB es la publicada por el Banco), activos y
pasivos de la posición de inversión internacional y PII neta sobre el
PIB en dólares. Reservas en meses de importaciones: reservas
internacionales netas sobre la doceava parte de las importaciones de
bienes de cuatro trimestres.

\textbf{Cuentas del Gobierno nacional central.} Las series del GNC son
de \textbf{caja} (empalme de series del DNP y del Ministerio de Hacienda
publicado por el Banco), no de causación. Los flujos mensuales se suman
por trimestre (solo trimestres completos), luego en cuatro trimestres, y
se dividen por el PIB nominal en pesos de los mismos cuatro trimestres:
\[\text{Z\%}_q=100\,\frac{\sum_{j=0}^{3}Z_{q-j}}{\sum_{j=0}^{3}PIB_{q-j}},\qquad Z\in\{\text{ingresos, gastos, intereses, balance, financiamiento}\}.\]
El balance primario es el balance total más los intereses,
\(BP=BT+INT\). Los intereses se expresan también por cada 100 pesos de
ingresos, \(100\,INT^{(4)}/I^{(4)}\). Financiamiento interno y externo
neto en \% del PIB, con la identidad \(F^{int}+F^{ext}=-BT\). La deuda
bruta del GNC es la razón anual publicada (serie 15328), comparada con
cinco años antes y con una referencia visual de 60\% del PIB. Si la
serie de intereses de caja tiene meses negativos en los últimos dos
años, el texto lo advierte (sección 9). Las series del sector público no
financiero se descargan pero no alimentan gráficos.

\hypertarget{comercio-exterior}{%
\subsubsection{5.12 Comercio exterior}\label{comercio-exterior}}

\textbf{Exportaciones e importaciones.} Exportaciones FOB y totales de
importaciones CIF (sistema comercial especial ampliado, incluye zonas
francas) del DANE con registros de la DIAN. El periodo de referencia
común es el último mes disponible en ambas series. Sumas de 12 meses,
variación frente a los 12 meses anteriores, balanza de 12 meses y
balanza por año calendario (el último año, hasta el último mes
publicado). Cobertura \(100\,X^{(12)}/M^{(12)}\).

\textbf{Qué se vende.} Valor de 12 meses y participación de no
tradicionales, petróleo y derivados, carbón, café y ferroníquel; peso
minero-energético \((P^{(12)}+C^{(12)})/X^{(12)}\). Se verifica que
tradicionales más no tradicionales reproduzcan el total con error menor
a 1\%.

\textbf{Precio y volumen.} Con los índices de precios de exportación e
importación en dólares del Banco (series 15361 y 15362): \(v\) =
variación anual del valor de 12 meses; \(p\) = variación anual del
promedio de 12 meses del índice de precios; volumen implícito
\(q=100\,[(1+v/100)/(1+p/100)-1]\). Términos de intercambio \(IPX/IPM\),
variación de 12 meses (Prebisch, 1950).

\textbf{Qué se compra (CUODE).} Variación del año corrido por grupo y
subgrupo de uso o destino económico y su participación (Cuadro A13); se
verifica que los grupos sumen el total con error menor a 2\%.
Composición anual entre bienes de consumo, materias primas y bienes de
capital y construcción (el último año es parcial).

\textbf{Socios.} Participación de cada destino en las exportaciones de
12 meses (nueve principales y ``resto''); de cada país de origen en las
importaciones de 12 meses (diez principales), con denominador igual a
las importaciones publicadas (sin zonas francas); balanza bilateral de
12 meses, exportaciones FOB menos importaciones CIF, para 11 socios;
participación de China y de Estados Unidos en las importaciones.

\textbf{Apertura y concentración.} Apertura de bienes: exportaciones más
importaciones de los últimos cuatro trimestres completos sobre el PIB en
dólares de los mismos trimestres. Número equivalente de destinos (16
destinos más ``resto'') y de productos (cinco grupos), \(1/H\) sobre
participaciones de 12 meses (Hausmann, Hwang y Rodrik, 2007; Hidalgo et
al., 2007).

\hypertarget{indicadores-y-datos}{%
\subsubsection{5.13 Indicadores y datos}\label{indicadores-y-datos}}

\textbf{Tabla de 16 indicadores.} Para crecimiento del PIB, brecha, ISE
(variación anual de la serie desestacionalizada), desempleo, inflación
total y de fondo, breakeven a 1 año, forward 5y5y, tasa de política,
tasa real ex ante, TES a 10 años, TRM, ITCR, COLCAP, cuenta corriente y
deuda del GNC: último dato, fecha, cambio (en pp o en \%, según la
naturaleza de la serie) frente a hace 365 días (o 91 días para
breakevens y TES a 10 años) y percentil histórico desde 2010, traducido
a palabras con los umbrales de la sección 6.

\textbf{Publicaciones oficiales.} Una nota resume el último dato
publicado por cada entidad (DANE, Banco de la República, BVC) a partir
de los datos descargados y enlaza a la página oficial; no copia
titulares de prensa. Cuando la fecha de publicación no está en el
archivo fuente se indica ``dato más reciente''.

\textbf{Estado de las fuentes y descargas.} La tabla de fuentes muestra
la última observación y el estado de cada fuente (sección 7). Los
archivos CSV se publican tal como se usan, para que cualquier cifra
pueda reproducirse.

\hypertarget{reglas-de-los-estados-y-del-veredicto-automuxe1tico}{%
\subsection{6. Reglas de los estados y del veredicto
automático}\label{reglas-de-los-estados-y-del-veredicto-automuxe1tico}}

\hypertarget{paruxe1metros-del-modelo}{%
\subsubsection{6.1 Parámetros del
modelo}\label{paruxe1metros-del-modelo}}

\small\begin{longtable}[]{@{}
  >{\raggedright\arraybackslash}p{(\columnwidth - 4\tabcolsep) * \real{0.3333}}
  >{\raggedright\arraybackslash}p{(\columnwidth - 4\tabcolsep) * \real{0.3333}}
  >{\raggedright\arraybackslash}p{(\columnwidth - 4\tabcolsep) * \real{0.3333}}@{}}
\toprule\noalign{}
\begin{minipage}[b]{\linewidth}\raggedright
Parámetro
\end{minipage} & \begin{minipage}[b]{\linewidth}\raggedright
Valor
\end{minipage} & \begin{minipage}[b]{\linewidth}\raggedright
Significado
\end{minipage} \\
\midrule\noalign{}
\endhead
\bottomrule\noalign{}
\endlastfoot
\texttt{NEUTRAL\_REAL} & \((2{,}7;\ 3{,}0)\) & Rango de la tasa real
neutral (\% real) \\
\texttt{UMBRAL\_BRECHA} & 0,5 & Brecha ``cerca del potencial'' si
\(\lvert g\rvert<0{,}5\%\) \\
\texttt{UMBRAL\_POSTURA} & 0,5 & Margen en pp alrededor del rango
neutral \\
\texttt{ANCLA\_EXPECTATIVAS} & 4,0 & Techo del rango meta para la
5y5y \\
Meta de inflación & \(3 \pm 1\) & Rango meta del Banco de la
República \\
\texttt{HP\_LAMBDA} & 1.600 / 129.600 & Suavizamiento trimestral /
mensual \\
\texttt{MIN\_OBS\_HP} & 20 trimestres & Historia mínima para la brecha
en tiempo real \\
\end{longtable}\normalsize

\hypertarget{veredicto-tuxe9cnico}{%
\subsubsection{6.2 Veredicto técnico}\label{veredicto-tuxe9cnico}}

El veredicto tiene un titular y cinco párrafos (actividad, precios,
política, expectativas, mercados); cada frase cita el dato que la
sustenta.

\small\begin{longtable}[]{@{}
  >{\raggedright\arraybackslash}p{(\columnwidth - 4\tabcolsep) * \real{0.3333}}
  >{\raggedright\arraybackslash}p{(\columnwidth - 4\tabcolsep) * \real{0.3333}}
  >{\raggedright\arraybackslash}p{(\columnwidth - 4\tabcolsep) * \real{0.3333}}@{}}
\toprule\noalign{}
\begin{minipage}[b]{\linewidth}\raggedright
Concepto
\end{minipage} & \begin{minipage}[b]{\linewidth}\raggedright
Regla
\end{minipage} & \begin{minipage}[b]{\linewidth}\raggedright
Resultado
\end{minipage} \\
\midrule\noalign{}
\endhead
\bottomrule\noalign{}
\endlastfoot
Fase & Reloj del ciclo con la brecha HP en tiempo real (sección 5.2) &
Expansión, desaceleración, contracción o recuperación \\
Posición frente al potencial & \(g\ge 0{,}5\); \(g\le -0{,}5\); en otro
caso & Por encima de; por debajo de; cerca de \\
Dirección de la inflación & \(\pi_t-\pi_{t-3}>0{,}1\) pp; \(<-0{,}1\)
pp; en otro caso & Acelerando; cediendo; estable \\
Rango meta & \(\pi_t>4\) o \(\pi_t<2\) & Fuera del rango meta; si no,
dentro \\
Postura monetaria & \(r^{ea}>3{,}0+0{,}5\); \(r^{ea}<2{,}7-0{,}5\); en
otro caso & Restrictiva; expansiva; neutral \\
Expectativas & \(\pi^{5y5y}>4\) & Desancladas; si no, ancladas \\
\end{longtable}\normalsize

\hypertarget{estados-para-no-especialistas-portada}{%
\subsubsection{6.3 Estados para no especialistas
(portada)}\label{estados-para-no-especialistas-portada}}

\small\begin{longtable}[]{@{}
  >{\raggedright\arraybackslash}p{(\columnwidth - 4\tabcolsep) * \real{0.3333}}
  >{\raggedright\arraybackslash}p{(\columnwidth - 4\tabcolsep) * \real{0.3333}}
  >{\raggedright\arraybackslash}p{(\columnwidth - 4\tabcolsep) * \real{0.3333}}@{}}
\toprule\noalign{}
\begin{minipage}[b]{\linewidth}\raggedright
Tarjeta
\end{minipage} & \begin{minipage}[b]{\linewidth}\raggedright
Regla
\end{minipage} & \begin{minipage}[b]{\linewidth}\raggedright
Estados
\end{minipage} \\
\midrule\noalign{}
\endhead
\bottomrule\noalign{}
\endlastfoot
Crecimiento & Con \(g^Y\) el PIB anual y \(g^*\) el crecimiento
potencial (HP de dos colas): \(g^Y<0\); \(g^Y\ge g^*+1\);
\(g^Y\ge g^*-1\); en otro caso & Contracción; fuerte; normal; lento \\
Inflación & \(2\le\pi\le4\); \(\pi<2\); \(\pi>5\); \(4<\pi\le5\) & En
meta; baja; alta; sobre la meta \\
Desempleo & Frente a hace 12 meses: \(TD<TD_{-12}-0{,}3\);
\(TD>TD_{-12}+0{,}3\); en otro caso & Mejora; empeora; estable \\
Tasa del Banco & Postura de la sección 6.2; si hay decisión anunciada
pendiente, ``sube'' o ``baja'' & Restrictiva, neutral, expansiva \\
Dólar & Variación de 12 meses de la TRM: \(<-5\%\); \(>5\%\); en otro
caso & Peso fuerte; peso débil; estable \\
Bolsa & Variación de 12 meses del COLCAP: \(>5\%\); \(<-5\%\); en otro
caso & Sube; baja; estable \\
Nivel histórico & Percentil \(<10\); \(<30\); \(\le 70\); \(\le 90\);
\(>90\) & Muy bajo; bajo; normal; alto; muy alto \\
\end{longtable}\normalsize

\hypertarget{seuxf1ales-de-color-en-las-tarjetas-de-cada-puxe1gina}{%
\subsubsection{6.4 Señales de color en las tarjetas de cada
página}\label{seuxf1ales-de-color-en-las-tarjetas-de-cada-puxe1gina}}

Las tarjetas de ``medidas'' de cada página llevan un tono (favorable o
de alerta) asignado por reglas fijas:

\small\begin{longtable}[]{@{}
  >{\raggedright\arraybackslash}p{(\columnwidth - 2\tabcolsep) * \real{0.5000}}
  >{\raggedright\arraybackslash}p{(\columnwidth - 2\tabcolsep) * \real{0.5000}}@{}}
\toprule\noalign{}
\begin{minipage}[b]{\linewidth}\raggedright
Página
\end{minipage} & \begin{minipage}[b]{\linewidth}\raggedright
Regla de alerta (o de tono favorable)
\end{minipage} \\
\midrule\noalign{}
\endhead
\bottomrule\noalign{}
\endlastfoot
Ciclo & Fase: expansión y recuperación favorables, desaceleración
alerta, contracción grave. Consenso favorable si al menos 3 de 5 métodos
dan brecha positiva. Amplitud ``generalizada'' si más de 6 de 12
sectores están sobre su tendencia. Empleo favorable si el desempleo está
bajo su tendencia \\
Capacidad & Alerta si la mediana de la brecha supera 0,5\%; sectores
favorables si más de 6 de 12 están sobre su tendencia; alerta si el
desempleo está bajo su tendencia (mercado apretado); subutilización
favorable si no supera la de un año antes; regiones favorables si al
menos 25 departamentos superan su nivel de 2019 \\
Crecimiento & Favorable si el crecimiento anual, el anualizado, el año
corrido o el crecimiento sin Gobierno es al menos 2\%; el texto atribuye
al sector público ``una parte grande'' del dato si el total supera al
crecimiento sin Gobierno en más de 0,5 pp; PIB por persona favorable si
crece; inversión favorable si la tasa supera su promedio de 2015 \\
Empleo & Favorable si suben los ocupados o la tasa de ocupación, si
bajan el desempleo, la informalidad o el desempleo juvenil frente a un
año antes, si los asalariados superan su participación de diciembre de
2019 y si la proporción de jóvenes que no estudian ni están ocupados no
supera su promedio de 2019; la brecha por sexo siempre en alerta \\
Inflación & Favorable si la inflación (total, servicios, no durables,
energéticos) está entre 2\% y 4\%; durables favorables si no superan
4\%; difusión en alerta si más de la mitad de las subclases supera
4\% \\
Tasas & Tasa real en alerta fuera de \([2{,}2;\,3{,}5]\); crédito total
en alerta si su crecimiento real es negativo \\
Curva TES & Pendiente en alerta si es menor a 0,5 pp; compensación a 10
años en alerta si supera 4\%; volatilidad en alerta si supera su
percentil 80 histórico \\
Mercados & Movimiento ``propio de Colombia'' si
\(\lvert v_{COP}-\bar v_{pares}\rvert>\max(3,\lvert\bar v_{pares}\rvert)\);
volatilidad en alerta sobre su percentil 80 \\
Empresas & Crédito a empresas en alerta si su crecimiento real es
negativo \\
Externo & Cuenta corriente en alerta si el déficit supera 4\% del PIB;
balance del Gobierno en alerta si el déficit supera 3\% del PIB; deuda
del GNC en alerta sobre 60\% del PIB \\
Comercio & Exportaciones favorables si crecen; balanza en alerta si es
deficitaria; bienes de capital favorables si crecen; balanza con China
en alerta si es deficitaria \\
\end{longtable}\normalsize

\hypertarget{reglas-de-texto-adicionales}{%
\subsubsection{6.5 Reglas de texto
adicionales}\label{reglas-de-texto-adicionales}}

\begin{itemize}
\tightlist
\item
  \emph{Curva:} forma normal si \(S>0{,}3\) pp, invertida si \(S<0\),
  plana en otro caso; movimiento de 12 meses según la sección 5.8.
\item
  \emph{Traspaso:} ``mayor en el interbancario y el crédito que en el
  ahorro'' si el cambio del CDT es menor que 60\% del menor de los
  cambios (en valor absoluto) del IBR a 3 meses y de la colocación;
  ``parecido'' en otro caso.
\item
  \emph{Posición ocupacional:} ``el empleo se formaliza por la vía de
  los contratos'' si los asalariados privados crecen más (en miles) que
  los trabajadores por cuenta propia.
\item
  \emph{Inflación por ingresos:} el texto indica qué grupo enfrenta la
  mayor inflación según el signo de la diferencia entre ingresos altos y
  pobres.
\item
  \emph{Deflactor:} el texto indica si el deflactor del PIB va por
  encima o por debajo del IPC del trimestre.
\end{itemize}

\hypertarget{control-de-calidad-y-automatizaciuxf3n}{%
\subsection{7. Control de calidad y
automatización}\label{control-de-calidad-y-automatizaciuxf3n}}

\hypertarget{controles-al-descargar}{%
\subsubsection{7.1 Controles al
descargar}\label{controles-al-descargar}}

Cada proceso de descarga valida sus datos antes de escribir y, si algo
no cuadra, conserva la versión anterior del archivo; una fuente caída
nunca borra a las demás.

\begin{itemize}
\tightlist
\item
  \textbf{Identidad.} Identificador y nombre de cada serie SUAMECA
  (sección 3.3); patrón y unicidad del anexo vigente del DANE;
  coincidencia de la fecha de publicación de los anexos real y nominal
  del PIB y del último trimestre con el nombre del archivo.
\item
  \textbf{Rangos y longitud mínima.} Rangos plausibles por serie
  (sección 3.2), historia mínima (por ejemplo, 1.000 datos diarios, 60
  mensuales, 40 trimestrales en \texttt{series\_banrep.csv}; al menos 60
  trimestres de PIB; 4.000 días de COLCAP; 200 meses de ISE y GEIH; 26
  semanas por acción), crecimientos fuera de control (PIB anual mayor a
  40\% o trimestral mayor a 30\%; sectores mayor a 80\%) y saltos
  diarios del COLCAP mayores a 25\%.
\item
  \textbf{Coherencia interna.} Departamentos que suman el total nacional
  (1\%); PIB del gasto igual al de producción (0,2\%); exportaciones
  tradicionales más no tradicionales iguales al total (1\%); grupos
  CUODE que suman el total (2\%); ponderaciones de divisiones del IPC
  que suman 100 (0,5) y de subclases (1,5); variaciones del IPC
  calculadas frente a las publicadas por el DANE (0,05 pp para el último
  mes y 0,10 pp para la historia mensual).
\item
  \textbf{Fechas.} Sin duplicados, sin fechas futuras (la meta de
  inflación del año siguiente se descarta), sin meses o trimestres
  faltantes en las series principales.
\item
  \textbf{Reintentos.} Reintentos con espera ante respuestas de límite
  de tasa del DANE (código 429) o fallas transitorias del servidor;
  pausas de cortesía entre series del Banco.
\end{itemize}

\hypertarget{validaciuxf3n-de-las-tablas-publicadas}{%
\subsubsection{7.2 Validación de las tablas
publicadas}\label{validaciuxf3n-de-las-tablas-publicadas}}

Después de descargar, un validador revisa las tablas y genera
\texttt{estado\_fuentes.csv}:

\begin{itemize}
\tightlist
\item
  Fuentes principales (PIB, IPC, TES, COLCAP): archivo presente, fechas
  válidas, ordenadas y sin futuros; PIB trimestral completo y
  variaciones recalculadas desde los niveles (tolerancia \(10^{-4}\));
  una sola vintage en la tabla del PIB; IPC mensual completo, sin
  variaciones faltantes, dentro de rangos de control y con el último mes
  contrastado con el DANE; TES en pesos completos y entre \(-5\%\) y
  40\%; COLCAP positivo, sin saltos diarios mayores a 25\% y con la base
  100 del 9 de febrero de 2009 recalculada.
\item
  Fuentes complementarias (ISE, GEIH, informalidad, series del Banco,
  canasta y acciones, sectores): meses consecutivos, identidad
  \(TD=100\,(1-TO/TGP)\) con tolerancia de 0,05 pp, variación anual del
  ISE recalculada, 12 agrupaciones por trimestre en sectores, tasas de
  informalidad entre 0 y 100, pesos de la canasta entre 80\% y 101\%.
\item
  \textbf{Estado de cada fuente.} ``Vigente'' si la antigüedad de la
  última observación (o de la última publicación, en PIB e IPC) no
  supera una tolerancia por fuente; ``rezagado'' si la supera;
  ``pendiente'' si aún no hay archivo.
\end{itemize}

\small\begin{longtable}[]{@{}ll@{}}
\toprule\noalign{}
Fuente & Tolerancia (días) \\
\midrule\noalign{}
\endhead
\bottomrule\noalign{}
\endlastfoot
PIB real & 130 desde la última publicación \\
IPC & 45 desde la última publicación \\
TES & 10 \\
COLCAP & 7 \\
ISE, informalidad & 110 \\
GEIH (desempleo) & 100 \\
Tasa de política, TRM & 10 \\
Inflación básica & 75 \\
ITCR & 100 \\
Cuenta corriente & 200 \\
Deuda del Gobierno & 800 \\
PIB por sectores & 230 \\
Canasta COLCAP, acciones & 10 \\
\end{longtable}\normalsize

Si el validador encuentra un error (no una advertencia), el proceso
termina con falla y el sitio no se publica.

\hypertarget{saltos-aislados}{%
\subsubsection{7.3 Saltos aislados}\label{saltos-aislados}}

En las tasas TES diarias, una observación se marca como salto aislado
cuando se aleja más de 2 pp de sus dos vecinos mientras los vecinos
difieren entre sí menos de 1 pp:
\[\lvert x_t-x_{t-1}\rvert>2,\quad \lvert x_t-x_{t+1}\rvert>2,\quad \lvert x_{t-1}-x_{t+1}\rvert<1.\]
Los saltos de las tasas en pesos se registran en
\texttt{alertas\_datos.csv}; en las seis series TES se conservan en el
CSV pero se excluyen de los cálculos derivados. Para las series
derivadas (breakevens a 1, 5 y 10 años, forward 5y5y y tasa real ex
ante), donde un salto de un día en una sola tasa UVR se amplifica, se
aplica el mismo criterio con umbrales de 1 pp y 0,5 pp.

\hypertarget{automatizaciuxf3n}{%
\subsubsection{7.4 Automatización}\label{automatizaciuxf3n}}

Un flujo de GitHub Actions actualiza los datos y publica el sitio
estático en GitHub Pages:

\begin{itemize}
\tightlist
\item
  \textbf{Actualización completa} cada noche a las 19:30, hora de
  Colombia: todas las fuentes, en secuencia y cada una en su propio
  proceso con un tiempo límite, de modo que la falla de una no detiene a
  las demás.
\item
  \textbf{Actualización rápida} cada hora, de 7:07 a 18:07 en días
  hábiles: curvas TES, COLCAP y series del Banco de la República
  (política, TRM, inflación básica, externas).
\item
  \textbf{Orden de cada corrida:} pruebas automáticas del código;
  descarga; aviso de fuentes con falla (no detiene la publicación: se
  usa el último dato válido); validación de tablas (un error detiene la
  publicación); pruebas con los datos del día; registro de los datos
  nuevos en el repositorio; construcción y publicación del sitio.
\item
  \textbf{Publicación.} La corrida nocturna publica siempre; la horaria
  solo si llegó un dato nuevo. Un cambio en el código reconstruye el
  sitio con los datos ya guardados, sin descargar.
\item
  \textbf{Trazabilidad.} El registro de cada corrida
  (\texttt{registro\_actualizacion.log}) y el estado de las fuentes se
  publican en el repositorio; cada actualización de datos queda como un
  cambio con fecha y modo (rápida o completa).
\end{itemize}

\hypertarget{ejercicio-de-robustez-fechado}{%
\subsection{8. Ejercicio de robustez
fechado}\label{ejercicio-de-robustez-fechado}}

\emph{Cálculo realizado el 3 de octubre de 2026 con los datos vigentes
en esa fecha (PIB hasta el T2 de 2026, publicado el 18 de agosto de
2026). Las cifras de esta sección cambian con cada publicación y
revisión; no son propiedades del método.}

\textbf{Brecha del producto.} Con 51 trimestres en los que los tres
métodos originales están disponibles (2013-T4 a 2026-T2):

\small\begin{longtable}[]{@{}
  >{\raggedright\arraybackslash}p{(\columnwidth - 2\tabcolsep) * \real{0.5000}}
  >{\raggedright\arraybackslash}p{(\columnwidth - 2\tabcolsep) * \real{0.5000}}@{}}
\toprule\noalign{}
\begin{minipage}[b]{\linewidth}\raggedright
Métrica
\end{minipage} & \begin{minipage}[b]{\linewidth}\raggedright
Valor
\end{minipage} \\
\midrule\noalign{}
\endhead
\bottomrule\noalign{}
\endlastfoot
Coincidencia de signo, HP en tiempo real frente a HP de dos colas &
78\% \\
Coincidencia de signo, HP en tiempo real frente a Hamilton & 65\% \\
Coincidencia de fase, HP en tiempo real frente a Hamilton (49
trimestres) & 43\% \\
Revisión media, tiempo real menos dos colas & \(-0{,}52\) pp (error
absoluto medio 0,96 pp; correlación 0,95) \\
Correlación entre los tres métodos sin 2020--2021 & 0,70 a 0,78 \\
\end{longtable}\normalsize

Lectura: el signo de la brecha es razonablemente robusto, pero la
\textbf{fase} depende del método. Por eso el sitio muestra el rango
entre métodos y el número de métodos con brecha positiva. En el último
trimestre la brecha en tiempo real coincide por construcción con la de
dos colas (ambas usan la misma muestra), de modo que la revisión futura
de ese dato solo se conocerá con trimestres adicionales.

\textbf{Ley de Okun.} Con 62 trimestres fuera de 2020--2021, la
pendiente estimada es \(\hat\beta=-0{,}12\) pp de desempleo bajo
tendencia por cada punto de brecha y la correlación es \(-0{,}26\): el
desempleo responde poco al ciclo.

\textbf{Inflación implícita a un año como predictor.} Promedio mensual
del breakeven a un año frente a la inflación observada doce meses
después, con la caminata aleatoria (inflación de hoy) como referencia:

\small\begin{longtable}[]{@{}
  >{\raggedright\arraybackslash}p{(\columnwidth - 8\tabcolsep) * \real{0.2000}}
  >{\raggedright\arraybackslash}p{(\columnwidth - 8\tabcolsep) * \real{0.2000}}
  >{\raggedright\arraybackslash}p{(\columnwidth - 8\tabcolsep) * \real{0.2000}}
  >{\raggedright\arraybackslash}p{(\columnwidth - 8\tabcolsep) * \real{0.2000}}
  >{\raggedright\arraybackslash}p{(\columnwidth - 8\tabcolsep) * \real{0.2000}}@{}}
\toprule\noalign{}
\begin{minipage}[b]{\linewidth}\raggedright
Periodo
\end{minipage} & \begin{minipage}[b]{\linewidth}\raggedright
Meses
\end{minipage} & \begin{minipage}[b]{\linewidth}\raggedright
Sesgo (observada − implícita)
\end{minipage} & \begin{minipage}[b]{\linewidth}\raggedright
RMSE implícita
\end{minipage} & \begin{minipage}[b]{\linewidth}\raggedright
RMSE caminata aleatoria
\end{minipage} \\
\midrule\noalign{}
\endhead
\bottomrule\noalign{}
\endlastfoot
2006--2019 & 168 & +0,48 pp & 1,70 pp & 2,08 pp \\
2020--2025 & 68 & +2,32 pp & 4,13 pp & 4,07 pp \\
Total & 236 & +1,01 pp & 2,64 pp & 2,80 pp \\
\end{longtable}\normalsize

Antes de 2020 la inflación implícita anticipó mejor que la caminata
aleatoria; el choque inflacionario de 2021--2023 no fue anticipado por
ninguno de los dos. El sitio presenta la implícita como \textbf{lo que
el mercado descuenta}, no como pronóstico.

\textbf{Saltos en series derivadas.} El filtro de la sección 7.3 excluye
110 observaciones diarias en total entre las cinco series derivadas.

\hypertarget{limitaciones-conocidas}{%
\subsection{9. Limitaciones conocidas}\label{limitaciones-conocidas}}

\begin{enumerate}
\def\labelenumi{\arabic{enumi}.}
\tightlist
\item
  \textbf{Una sola vintage.} Cada corrida reemplaza la historia con la
  última publicación del DANE y del Banco. Las cifras históricas son las
  revisadas, no las que se conocían en cada fecha. La brecha ``en tiempo
  real'' elimina la información futura del filtro, pero no las
  revisiones de los datos; la fecha de publicación registrada es la del
  anexo vigente, no la de primera publicación de cada periodo.
\item
  \textbf{Datos provisionales.} Los últimos trimestres del PIB, del PIB
  por el gasto y de las cuentas departamentales son provisionales y se
  revisan; las cuentas departamentales llegan con más de un año de
  rezago.
\item
  \textbf{Sensibilidad al método.} La brecha y la fase del ciclo cambian
  con el método (sección 8). Cuatro de los cinco métodos usan toda la
  muestra y su último dato se revisará. El tratamiento de la pandemia
  (interpolación de 2020-T2 a 2021-T2) es una decisión de modelado.
\item
  \textbf{Sesgo de fin de muestra del fechado de giros.} El promedio
  centrado y la ventana de \(\pm 6\) meses usan meses posteriores; el
  último giro puede moverse o desaparecer con datos nuevos.
\item
  \textbf{Tasa neutral externa.} El rango de 2,7\% a 3,0\% es una
  estimación con incertidumbre considerable y se actualiza a mano; la
  clasificación de la postura depende de él.
\item
  \textbf{Primas en los breakevens.} La inflación implícita incluye
  primas por riesgo inflacionario y por liquidez (los TES UVR son menos
  líquidos) y el rezago de indexación de la UVR; no es una expectativa
  pura (Gürkaynak, Sack y Wright, 2010). Las primas pueden tener
  cualquier signo, de modo que la implícita no es necesariamente una
  cota superior.
\item
  \textbf{Curva estimada.} Las tasas cero cupón son salidas de un modelo
  (Nelson-Siegel) estimado por el Banco, no precios negociados; los
  plazos de 1, 5 y 10 años no describen toda la curva.
\item
  \textbf{Índice de precios frente a retorno total.} El COLCAP y los
  índices equiponderado y de 7 Magníficas son índices de precios, sin
  dividendos; no miden el retorno total del inversionista.
\item
  \textbf{Fuentes auxiliares y sesgo de supervivencia.} La bolsa por
  dentro usa pesos de un fondo (no los oficiales) y precios de un
  proveedor comercial sin ajuste por eventos corporativos; aplicar la
  canasta vigente hacia atrás omite las empresas que salieron del índice
  (sesgo de supervivencia) y el empalme de tickers es una aproximación.
\item
  \textbf{10.000 empresas.} La lista cambia cada año, excluye bancos y
  aseguradoras y viene en billones con dos decimales; las sumas pueden
  diferir levemente del informe oficial y las variaciones no comparan
  las mismas empresas. Las regiones corresponden al domicilio, no al
  lugar de producción.
\item
  \textbf{Registro mercantil.} Las cancelaciones incluyen depuraciones
  administrativas del registro y presentan saltos; el último mes puede
  estar incompleto aun después del filtro.
\item
  \textbf{FOB frente a CIF.} Las balanzas comercial y bilateral restan
  importaciones CIF (con fletes y seguros) de exportaciones FOB, lo que
  exagera el déficit frente a la balanza de pagos (FOB/FOB). La
  participación por país de origen usa como denominador las
  importaciones publicadas sin zonas francas.
\item
  \textbf{Volumen implícito.} El volumen de comercio es residual: hereda
  cualquier diferencia de cobertura entre los valores del DANE y los
  índices de precios del Banco.
\item
  \textbf{Cuentas fiscales en caja.} Las series del GNC son de caja, sin
  causaciones; la serie de intereses puede tener meses negativos, por lo
  que los intereses de 12 meses y la razón intereses/ingresos pueden
  subestimar el costo de la deuda que el Ministerio de Hacienda reporta
  en causación.
\item
  \textbf{Razones al PIB en dólares.} Dependen de la TRM promedio usada
  para convertir el PIB; un movimiento cambiario fuerte altera la razón
  sin cambio en el numerador. Algunas razones (cuenta corriente
  histórica, deuda externa) son las publicadas por el Banco con su
  propio PIB de referencia y pueden diferir de las calculadas en el
  sitio.
\item
  \textbf{Aportes con volúmenes encadenados.} Los aportes sectoriales y
  de la demanda no suman exactamente el total; el crecimiento ``sin
  Gobierno'' o ``sin minería'' es una aproximación con pesos corrientes.
\item
  \textbf{Mezcla de vintages entre fuentes.} El PIB, el IPC, la GEIH y
  las series del Banco se actualizan en días distintos; una misma página
  puede combinar datos de meses diferentes, lo que se indica con la
  fecha de cada cifra.
\item
  \textbf{Identidad parcial de algunas series.} Las series de cuentas
  externas, fiscales y de empresas del Banco se validan por
  identificador y longitud, sin texto esperado en el nombre; una
  reasignación de identificador que conserve la longitud no se
  detectaría automáticamente.
\item
  \textbf{Informalidad y GEIH.} La serie de informalidad empieza en 2021
  (marco muestral 2018) y no es comparable con la anterior; las series
  de detalle no están desestacionalizadas.
\item
  \textbf{Registro manual.} La decisión anunciada de la Junta se
  registra a mano (sección 3.5); un error de digitación afectaría la
  tarjeta de la tasa hasta que la decisión entre en vigencia.
\end{enumerate}

\hypertarget{referencias}{%
\subsection{10. Referencias}\label{referencias}}

Arango, L. E., Melo, L. F. y Vásquez, D. M. (2002). \emph{Estimación de
la estructura a plazo de las tasas de interés en Colombia} (Borradores
de Economía n.º 196). Banco de la República.

Banco de la República. (2025). \emph{Informe de Política Monetaria}
(enero, abril y octubre de 2025) y estimaciones de la tasa de interés
neutral del equipo técnico. Banco de la República.

Banco de la República. (s. f.). \emph{Definiciones de la tasa de
política monetaria, IBR, TIB, DTF, CDT y tasas de colocación} (catálogo
de series estadísticas). Banco de la República.

Banco de la República. (s. f.). \emph{Mecanismos de transmisión de la
política monetaria en Colombia}. Banco de la República.

Banco de la República. (s. f.). \emph{Metodología de cálculo del Índice
de Tasa de Cambio Real (ITCR) de Colombia}. Banco de la República.

Banco de la República. (s. f.). \emph{Reporte de Estabilidad Financiera}
(semestral). Banco de la República.

Banco de la República, Junta Directiva. (2018). \emph{Resolución Externa
1 de 2018} y Circular Reglamentaria DOAM-146. Banco de la República.

Baumol, W. J. (1967). Macroeconomics of unbalanced growth: The anatomy
of urban crisis. \emph{American Economic Review, 57}(3), 415--426.

Baxter, M. y King, R. G. (1999). Measuring business cycles: Approximate
band-pass filters for economic time series. \emph{Review of Economics
and Statistics, 81}(4), 575--593.

Bell, D. N. F. y Blanchflower, D. G. (2011). Young people and the Great
Recession. \emph{Oxford Review of Economic Policy, 27}(2), 241--267.

Bernanke, B. S. y Gertler, M. (1995). Inside the black box: The credit
channel of monetary policy transmission. \emph{Journal of Economic
Perspectives, 9}(4), 27--48.

Bernanke, B. S., Gertler, M. y Gilchrist, S. (1999). The financial
accelerator in a quantitative business cycle framework. En J. B. Taylor
y M. Woodford (Eds.), \emph{Handbook of Macroeconomics} (Vol. 1C,
pp.~1341--1393). Elsevier.

Betancourt, R., Vargas, H. y Rodríguez, N. (2008). Interest rate
pass-through in Colombia: A micro-banking perspective. \emph{Cuadernos
de Economía, 45}(131), 29--58.

Beveridge, S. y Nelson, C. R. (1981). A new approach to decomposition of
economic time series into permanent and transitory components with
particular attention to measurement of the ``business cycle''.
\emph{Journal of Monetary Economics, 7}(2), 151--174.

Blanchard, O. (2019). Public debt and low interest rates. \emph{American
Economic Review, 109}(4), 1197--1229.

Blanchard, O., Cerutti, E. y Summers, L. (2015). \emph{Inflation and
activity: Two explorations and their monetary policy implications} (IMF
Working Paper WP/15/230). Fondo Monetario Internacional.

Bry, G. y Boschan, C. (1971). \emph{Cyclical analysis of time series:
Selected procedures and computer programs}. National Bureau of Economic
Research.

Bryan, M. F. y Cecchetti, S. G. (1994). Measuring core inflation. En N.
G. Mankiw (Ed.), \emph{Monetary Policy} (pp.~195--215). University of
Chicago Press.

Burns, A. F. y Mitchell, W. C. (1946). \emph{Measuring business cycles}.
National Bureau of Economic Research.

Calvo, G. A., Leiderman, L. y Reinhart, C. M. (1993). Capital inflows
and real exchange rate appreciation in Latin America: The role of
external factors. \emph{IMF Staff Papers, 40}(1), 108--151.

Cecchetti, S. G. (1997). Measuring short-run inflation for central
bankers. \emph{Federal Reserve Bank of St.~Louis Review, 79}(3),
143--155.

Cerra, V. y Saxena, S. C. (2008). Growth dynamics: The myth of economic
recovery. \emph{American Economic Review, 98}(1), 439--457.

Chavarro, X., Cristiano, D., Gómez, J. E., González, E. y Huertas, C.
(2015). \emph{Evaluación de la transmisión de la tasa de interés de
referencia a las tasas de interés del sistema financiero} (Borradores de
Economía n.º 874). Banco de la República.

Chen, Y.-C. y Rogoff, K. (2003). Commodity currencies. \emph{Journal of
International Economics, 60}(1), 133--160.

Christiano, L. J. y Fitzgerald, T. J. (2003). The band pass filter.
\emph{International Economic Review, 44}(2), 435--465.

Clarida, R., Galí, J. y Gertler, M. (1999). The science of monetary
policy: A New Keynesian perspective. \emph{Journal of Economic
Literature, 37}(4), 1661--1707.

Congreso de la República de Colombia. (2011). \emph{Ley 1473 de 2011,
por medio de la cual se establece una regla fiscal y se dictan otras
disposiciones}.

Congreso de la República de Colombia. (2013). \emph{Ley Estatutaria 1622
de 2013, por medio de la cual se expide el estatuto de ciudadanía
juvenil}.

Congreso de la República de Colombia. (2021). \emph{Ley 2155 de 2021,
por medio de la cual se expide la Ley de Inversión Social} (modifica la
regla fiscal y crea el Comité Autónomo de la Regla Fiscal).

DANE. (2019). \emph{Índice de Precios al Consumidor, base diciembre de
2018: metodología y canasta}. Departamento Administrativo Nacional de
Estadística.

DANE. (2021). \emph{Gran Encuesta Integrada de Hogares: metodología y
marco conceptual}. Departamento Administrativo Nacional de Estadística.

DANE. (s. f.). \emph{Cuentas nacionales departamentales, base 2015:
metodología (CD-01)}. Departamento Administrativo Nacional de
Estadística.

DANE. (s. f.). \emph{Cuentas nacionales trimestrales: metodología
(enfoque de la producción, índices encadenados)}. Departamento
Administrativo Nacional de Estadística.

DANE. (s. f.). \emph{Estadísticas de comercio internacional:
exportaciones (FOB) e importaciones (CIF) con base en registros
administrativos de la DIAN}. Departamento Administrativo Nacional de
Estadística.

DANE. (s. f.). \emph{Proyecciones y retroproyecciones de población,
Censo Nacional de Población y Vivienda 2018 (actualización 2025)}.
Departamento Administrativo Nacional de Estadística.

Dornbusch, R. (1976). Expectations and exchange rate dynamics.
\emph{Journal of Political Economy, 84}(6), 1161--1176.

Espinosa, J. A., Melo, L. F. y Moreno, J. F. (2014). \emph{Estimación de
la prima por vencimiento de los TES en pesos del gobierno colombiano}
(Borradores de Economía n.º 854). Banco de la República.

Estrella, A. y Hardouvelis, G. A. (1991). The term structure as a
predictor of real economic activity. \emph{Journal of Finance, 46}(2),
555--576.

Fisher, I. (1930). \emph{The theory of interest}. Macmillan.

Fondo Monetario Internacional. (2009). \emph{Manual de balanza de pagos
y posición de inversión internacional} (6.ª ed., MBP6). FMI.

Fondo Monetario Internacional, Organización Internacional del Trabajo,
OCDE, Eurostat, Comisión Económica de las Naciones Unidas para Europa y
Banco Mundial. (2004). \emph{Consumer price index manual: Theory and
practice}. FMI.

Gabaix, X. (2011). The granular origins of aggregate fluctuations.
\emph{Econometrica, 79}(3), 733--772.

Goldin, C. (2014). A grand gender convergence: Its last chapter.
\emph{American Economic Review, 104}(4), 1091--1119.

Gürkaynak, R. S., Sack, B. y Wright, J. H. (2010). The TIPS yield curve
and inflation compensation. \emph{American Economic Journal:
Macroeconomics, 2}(1), 70--92.

Hamilton, J. D. (2018). Why you should never use the Hodrick-Prescott
filter. \emph{Review of Economics and Statistics, 100}(5), 831--843.

Hausmann, R., Hwang, J. y Rodrik, D. (2007). What you export matters.
\emph{Journal of Economic Growth, 12}(1), 1--25.

Herfindahl, O. C. (1950). \emph{Concentration in the steel industry}
(tesis doctoral). Columbia University.

Hidalgo, C. A., Klinger, B., Barabási, A.-L. y Hausmann, R. (2007). The
product space conditions the development of nations. \emph{Science,
317}(5837), 482--487.

Hirschman, A. O. (1964). The paternity of an index. \emph{American
Economic Review, 54}(5), 761.

Hodrick, R. J. y Prescott, E. C. (1997). Postwar U.S. business cycles:
An empirical investigation. \emph{Journal of Money, Credit and Banking,
29}(1), 1--16.

Hsieh, C.-T. y Klenow, P. J. (2009). Misallocation and manufacturing TFP
in China and India. \emph{Quarterly Journal of Economics, 124}(4),
1403--1448.

Jaravel, X. (2021). Inflation inequality: Measurement, causes, and
policy implications. \emph{Annual Review of Economics, 13}, 599--629.

Kohli, U. (2004). Real GDP, real domestic income, and terms-of-trade
changes. \emph{Journal of International Economics, 62}(1), 83--106.

Lane, P. R. y Milesi-Ferretti, G. M. (2007). The external wealth of
nations mark II: Revised and extended estimates of foreign assets and
liabilities, 1970--2004. \emph{Journal of International Economics,
73}(2), 223--250.

Laubach, T. y Williams, J. C. (2003). Measuring the natural rate of
interest. \emph{Review of Economics and Statistics, 85}(4), 1063--1070.

Litterman, R. y Scheinkman, J. (1991). Common factors affecting bond
returns. \emph{Journal of Fixed Income, 1}(1), 54--61.

Maloney, W. F. (2004). Informality revisited. \emph{World Development,
32}(7), 1159--1178.

Meese, R. A. y Rogoff, K. (1983). Empirical exchange rate models of the
seventies: Do they fit out of sample? \emph{Journal of International
Economics, 14}(1--2), 3--24.

Melitz, M. J. (2003). The impact of trade on intra-industry
reallocations and aggregate industry productivity. \emph{Econometrica,
71}(6), 1695--1725.

Modigliani, F. y Miller, M. H. (1958). The cost of capital, corporation
finance and the theory of investment. \emph{American Economic Review,
48}(3), 261--297.

MSCI. (2021). \emph{MSCI COLCAP Index methodology}. MSCI Inc.

Naciones Unidas. (1999). \emph{Classification of individual consumption
according to purpose (COICOP)}. Naciones Unidas.

Naciones Unidas, Comisión Europea, FMI, OCDE y Banco Mundial. (2009).
\emph{System of National Accounts 2008}. Naciones Unidas.

Nelson, C. R. y Siegel, A. F. (1987). Parsimonious modeling of yield
curves. \emph{Journal of Business, 60}(4), 473--489.

Obstfeld, M. y Rogoff, K. (1995). The intertemporal approach to the
current account. En G. M. Grossman y K. Rogoff (Eds.), \emph{Handbook of
International Economics} (Vol. 3, pp.~1731--1799). Elsevier.

OECD. (2001). \emph{Measuring productivity: OECD manual}. OECD
Publishing.

OECD. (s. f.). \emph{Business Cycle Clock} (metodología del sistema de
indicadores adelantados compuestos). OECD.

Okun, A. M. (1962). Potential GNP: Its measurement and significance. En
\emph{Proceedings of the Business and Economic Statistics Section}
(pp.~98--104). American Statistical Association.

Organización Internacional del Trabajo. (1993). \emph{Resolución sobre
la Clasificación Internacional de la Situación en el Empleo (CISE-93)}.
15.ª Conferencia Internacional de Estadísticos del Trabajo.

Organización Internacional del Trabajo. (2013). \emph{Resolución sobre
las estadísticas del trabajo, la ocupación y la subutilización de la
fuerza de trabajo}. 19.ª Conferencia Internacional de Estadísticos del
Trabajo.

Orphanides, A. y van Norden, S. (2002). The unreliability of output-gap
estimates in real time. \emph{Review of Economics and Statistics,
84}(4), 569--583.

Prebisch, R. (1950). \emph{The economic development of Latin America and
its principal problems}. Naciones Unidas, CEPAL.

Rajan, R. G. y Zingales, L. (1998). Financial dependence and growth.
\emph{American Economic Review, 88}(3), 559--586.

Ravn, M. O. y Uhlig, H. (2002). On adjusting the Hodrick-Prescott filter
for the frequency of observations. \emph{Review of Economics and
Statistics, 84}(2), 371--376.

Reinhart, C. M. y Rogoff, K. S. (2009). \emph{This time is different:
Eight centuries of financial folly}. Princeton University Press.

Rey, H. (2013). Dilemma not trilemma: The global financial cycle and
monetary policy independence. En \emph{Proceedings of the Jackson Hole
Economic Policy Symposium} (pp.~285--333). Federal Reserve Bank of
Kansas City.

Rogoff, K. (1996). The purchasing power parity puzzle. \emph{Journal of
Economic Literature, 34}(2), 647--668.

Superintendencia de Sociedades. (2026). \emph{Informe de las 10.000
empresas más grandes del país} (cifras a diciembre de 2025).
Supersociedades.

Superintendencia Financiera de Colombia. (s. f.). \emph{Tasa de cambio
representativa del mercado: antecedentes normativos y metodología}.
Superfinanciera.

Taylor, J. B. (1993). Discretion versus policy rules in practice.
\emph{Carnegie-Rochester Conference Series on Public Policy, 39},
195--214.
